% =====================================================================
%  CHƯƠNG 2. ĐẶC TẢ YÊU CẦU BÀI TOÁN
%  Nội dung chia thành các tệp con trong chapters/c2/ để dễ cộng tác.
% =====================================================================
\chapter{ĐẶC TẢ YÊU CẦU BÀI TOÁN}\label{chap:dac-ta}

Chương này đặc tả yêu cầu của ứng dụng “Đi chợ tiện lợi” dưới dạng mô hình use case theo chuẩn UML~\cite{uml251}, kết hợp với các yêu cầu phi chức năng được tổ chức theo mô hình chất lượng phần mềm ISO/IEC 25010~\cite{iso25010}. Cách tiếp cận này tuân theo các khuyến nghị về kỹ nghệ yêu cầu trong ISO/IEC/IEEE 29148~\cite{iso29148}: mỗi yêu cầu được định danh, có thể kiểm chứng và truy vết được về nhu cầu nghiệp vụ đã phân tích ở Chương~\ref{chap:khao-sat}. Nội dung chương gồm ba phần. Phần thứ nhất xác định các tác nhân, các use case và quan hệ giữa chúng. Phần thứ hai trình bày biểu đồ use case tổng quan, các biểu đồ phân rã mức 2 cùng đặc tả chi tiết của các use case chính; các use case còn lại được đặc tả trong Phụ lục~\ref{pl:dac-ta-bo-sung}. Phần thứ ba xác định các yêu cầu phi chức năng.

Trong các bảng đặc tả, nhóm sử dụng thống nhất bốn loại mã tham chiếu: mã use case dạng \uc{UCxx}, mã quy tắc nghiệp vụ dạng \br{BRxx} đã định nghĩa trong Bảng~\ref{tab:quy-tac-nghiep-vu}, mã thông điệp hệ thống dạng \msg{MSGxx} được liệt kê trong Phụ lục~\ref{pl:thong-diep}, và mã yêu cầu phi chức năng dạng \nfr{NFRxx} được định nghĩa ở mục~\ref{sec:phi-chuc-nang}.

% ---------------------------------------------------------------------
%  2.1. Giới thiệu chung — tác nhân, use case, quan hệ
% ---------------------------------------------------------------------
\section{Giới thiệu chung}\label{sec:gioi-thieu-chung}

\subsection{Xác định các tác nhân của hệ thống}\label{ssec:tac-nhan}

Trong UML, tác nhân là vai trò do con người hoặc hệ thống bên ngoài đảm nhận khi tương tác với hệ thống đang xét~\cite{cockburn}. Từ sơ đồ ngữ cảnh ở Hình~\ref{fig:ngu-canh}, nhóm xác định chín tác nhân thuộc hai loại. Năm tác nhân chính chủ động khởi tạo use case để đạt mục tiêu: Khách, Người dùng, Thành viên nhóm, Trưởng nhóm và Quản trị viên. Bốn tác nhân phụ tham gia theo yêu cầu hoặc kích hoạt hệ thống theo lịch: Bộ lập lịch, Dịch vụ thông báo đẩy, Dịch vụ thư điện tử và Cơ sở dữ liệu sản phẩm mở. Bộ lập lịch được mô hình hóa riêng để thể hiện tác vụ tự động gửi thông báo khi thực phẩm sắp hết hạn. Mô tả các tác nhân được trình bày trong Bảng~\ref{tab:tac-nhan}.

{\small
\begin{longtable}{|C{0.9cm}|L{3.3cm}|L{\dimexpr\textwidth-6\tabcolsep-4\arrayrulewidth-0.9cm-3.3cm\relax}|}
\caption{Danh sách các tác nhân của hệ thống}\label{tab:tac-nhan}\\
\hline
\hd{STT} & \hd{Tên tác nhân} & \hd{Mô tả tác nhân}\\ \hline
\endfirsthead
\multicolumn{3}{l}{\footnotesize\itshape Bảng \thetable\ (tiếp theo)}\\ \hline
\hd{STT} & \hd{Tên tác nhân} & \hd{Mô tả tác nhân}\\ \hline
\endhead
1 & Khách & Người chưa có tài khoản hoặc chưa đăng nhập. Khách chỉ được đăng ký tài khoản, xác thực tài khoản, đăng nhập và khôi phục mật khẩu.\\ \hline
2 & Người dùng & Người đã đăng nhập thành công. Người dùng làm việc trong không gian hiện hành (không gian cá nhân hoặc nhóm gia đình) với đầy đủ chức năng quản lý danh sách mua sắm, tủ lạnh, công thức, kế hoạch bữa ăn, tìm kiếm và thống kê; có thể tạo nhóm gia đình mới hoặc tham gia nhóm có sẵn.\\ \hline
3 & Thành viên nhóm & Người dùng đã tham gia một nhóm gia đình. Ngoài các quyền của Người dùng, thành viên được xem thông tin nhóm, đóng góp vào dữ liệu chung, nhận nhiệm vụ mua sắm được phân công và có thể rời nhóm.\\ \hline
4 & Trưởng nhóm & Thành viên giữ quyền điều phối của nhóm, là người tạo nhóm hoặc người được chuyển quyền. Trưởng nhóm mời và xóa thành viên, chuyển quyền trưởng nhóm và phân công việc mua sắm, bảo đảm không bỏ sót hay mua trùng mặt hàng.\\ \hline
5 & Quản trị viên & Người dùng mang vai trò quản trị của hệ thống. Quản trị viên quản lý tài khoản người dùng, danh mục thực phẩm, loại mặt hàng, đơn vị đo lường, công thức mẫu và các tham số cấu hình, nhưng không truy cập nội dung dữ liệu của các nhóm (\br{BR23}).\\ \hline
6 & Bộ lập lịch & Tác nhân thời gian thuộc hệ thống, kích hoạt các tác vụ định kỳ như cập nhật trạng thái hạn dùng và gửi nhắc nhở thực phẩm sắp hết hạn vào khung giờ đã cấu hình.\\ \hline
7 & Dịch vụ thông báo đẩy & Hệ thống bên ngoài (Firebase Cloud Messaging cho Android, Apple Push Notification service cho iOS~\cite{fcm}) nhận yêu cầu từ máy chủ và chuyển thông báo tới thiết bị di động của người dùng.\\ \hline
8 & Dịch vụ thư điện tử & Hệ thống bên ngoài gửi thư chứa mã OTP khi đăng ký, khi khôi phục mật khẩu và các thư thông báo liên quan đến tài khoản.\\ \hline
9 & Cơ sở dữ liệu sản phẩm mở & Hệ thống bên ngoài cung cấp thông tin sản phẩm đóng gói theo mã vạch~\cite{openfoodfacts}, được tra cứu khi mã vạch chưa có trong danh mục nội bộ.\\ \hline
\end{longtable}
}

\subsection{Quan hệ giữa các tác nhân}\label{ssec:quan-he-tac-nhan}

Hình~\ref{fig:quan-he-tac-nhan} thể hiện hai nhóm tác nhân chính và phụ. Trong nhóm tác nhân chính, Thành viên nhóm là trường hợp chuyên biệt của Người dùng, còn Trưởng nhóm là trường hợp chuyên biệt của Thành viên nhóm. Theo quan hệ tổng quát hóa trong UML, tác nhân con kế thừa các use case của tác nhân cha; vì vậy, Trưởng nhóm kế thừa các use case của cả Thành viên nhóm lẫn Người dùng mà không cần lặp lại liên kết trên biểu đồ. Quản trị viên cũng kế thừa từ Người dùng để sử dụng các use case tài khoản như đăng xuất, đổi mật khẩu và cập nhật hồ sơ.

\diagram[0.95\textwidth]{c2-01-quan-he-tac-nhan}{Biểu đồ quan hệ giữa các tác nhân}{fig:quan-he-tac-nhan}

Khách và Người dùng không có quan hệ kế thừa; đây là hai vai trò mà một người có thể lần lượt đảm nhận sau khi đăng ký, xác thực hoặc đăng nhập. Tương tự, vai trò Thành viên nhóm và Trưởng nhóm phụ thuộc vào trạng thái tham gia nhóm: khi trưởng nhóm chuyển quyền, người nhận trở thành Trưởng nhóm và người chuyển quyền trở lại vai trò Thành viên nhóm. Các tác nhân phụ hoạt động độc lập, không có quan hệ tổng quát hóa với nhau hoặc với tác nhân chính.

\subsection{Xác định các ca sử dụng}\label{ssec:danh-sach-uc}

Từ các nghiệp vụ con ở mục~\ref{sec:thong-tin-nghiep-vu}, nhóm xác định 53 use case và tổ chức thành chín gói chức năng dựa trên tám phân hệ ở Hình~\ref{fig:phan-ra-chuc-nang}; phân hệ tài khoản và nhóm gia đình được tách thành hai gói. Mỗi use case có mã định danh để tham chiếu trong báo cáo. Độ phức tạp được chia thành ba mức, căn cứ vào số bước tương tác, số nhánh và mức độ phụ thuộc dữ liệu: \emph{Đơn giản} khi chủ yếu đọc hoặc ghi một thực thể với ít ràng buộc; \emph{Trung bình} khi có nhiều điều kiện kiểm tra hoặc tương tác với một phân hệ khác; \emph{Phức tạp} khi có thuật toán, giao dịch trên nhiều thực thể hoặc tương tác với hệ thống bên ngoài. Bảng~\ref{tab:danh-sach-uc} liệt kê các use case; ký hiệu (*) đánh dấu chức năng mở rộng.

{\small
\setlength{\LTleft}{0pt}
\begin{longtable}{|C{0.8cm}|C{1.25cm}|L{3.2cm}|L{\dimexpr\textwidth-12\tabcolsep-7\arrayrulewidth-0.8cm-1.25cm-3.2cm-2.6cm-1.85cm\relax}|L{2.6cm}|C{1.85cm}|}
\caption{Danh sách các use case của hệ thống}\label{tab:danh-sach-uc}\\
\hline
\hd{STT} & \hd{Mã} & \hd{Tên use case} & \hd{Mô tả} & \hd{Tác nhân} & \hd{Độ phức tạp}\\ \hline
\endfirsthead
\multicolumn{6}{l}{\footnotesize\itshape Bảng \thetable\ (tiếp theo)}\\ \hline
\hd{STT} & \hd{Mã} & \hd{Tên use case} & \hd{Mô tả} & \hd{Tác nhân} & \hd{Độ phức tạp}\\ \hline
\endhead
\multicolumn{6}{|l|}{\hs{Gói 1. Xác thực và quản lý tài khoản}}\\ \hline
1 & UC01 & Đăng ký tài khoản & Tạo tài khoản mới bằng họ tên, email và mật khẩu & Khách & Trung bình\\ \hline
2 & UC02 & Xác thực mã OTP & Xác minh quyền sở hữu email bằng mã gửi qua thư & Khách, Dịch vụ thư điện tử & Trung bình\\ \hline
3 & UC03 & Đăng nhập & Đăng nhập bằng email và mật khẩu để bắt đầu phiên làm việc & Khách & Đơn giản\\ \hline
4 & UC04 & Khôi phục mật khẩu & Đặt lại mật khẩu khi quên, thông qua mã OTP & Khách & Trung bình\\ \hline
5 & UC05 & Đăng xuất & Kết thúc phiên, thu hồi mã làm mới của thiết bị & Người dùng & Đơn giản\\ \hline
6 & UC06 & Cập nhật hồ sơ cá nhân & Sửa họ tên, ảnh đại diện, ngôn ngữ hiển thị & Người dùng & Đơn giản\\ \hline
7 & UC07 & Đổi mật khẩu & Đổi mật khẩu khi đã đăng nhập, thu hồi các phiên khác & Người dùng & Đơn giản\\ \hline
8 & UC08 & Thiết lập thông báo nhắc nhở (*) & Chọn ngưỡng nhắc 1–7 ngày, bật tắt từng loại thông báo & Người dùng & Đơn giản\\ \hline
\multicolumn{6}{|l|}{\hs{Gói 2. Quản lý nhóm gia đình}}\\ \hline
9 & UC09 & Tạo nhóm gia đình & Tạo nhóm mới, người tạo trở thành trưởng nhóm & Người dùng & Trung bình\\ \hline
10 & UC10 & Mời thành viên vào nhóm (*) & Sinh mã mời, mã QR có thời hạn để chia sẻ; thu hồi mã & Trưởng nhóm & Trung bình\\ \hline
11 & UC11 & Tham gia nhóm gia đình (*) & Gia nhập nhóm bằng mã mời hoặc quét mã QR & Người dùng & Trung bình\\ \hline
12 & UC12 & Xem thông tin nhóm và thành viên & Xem tên nhóm, danh sách thành viên và vai trò & Thành viên nhóm & Đơn giản\\ \hline
13 & UC13 & Xóa thành viên khỏi nhóm & Loại một thành viên ra khỏi nhóm & Trưởng nhóm & Đơn giản\\ \hline
14 & UC14 & Chuyển quyền trưởng nhóm & Trao vai trò trưởng nhóm cho một thành viên khác & Trưởng nhóm & Trung bình\\ \hline
15 & UC15 & Rời nhóm gia đình & Rời nhóm, trở về không gian cá nhân & Thành viên nhóm & Trung bình\\ \hline
\multicolumn{6}{|l|}{\hs{Gói 3. Quản lý danh sách mua sắm}}\\ \hline
16 & UC16 & Tạo danh sách mua sắm & Tạo danh sách gắn với ngày dự kiến mua & Người dùng & Trung bình\\ \hline
17 & UC17 & Cập nhật mặt hàng trong danh sách & Thêm, sửa, xóa mặt hàng; gộp mặt hàng trùng & Người dùng & Trung bình\\ \hline
18 & UC18 & Sao chép danh sách mua sắm đã có & Tạo danh sách mới từ một danh sách cũ để tái sử dụng & Người dùng & Đơn giản\\ \hline
19 & UC19 & Phân công người mua & Giao từng mặt hàng cho một thành viên, gửi thông báo & Trưởng nhóm, Dịch vụ thông báo đẩy & Phức tạp\\ \hline
20 & UC20 & Đánh dấu mặt hàng đã mua & Ghi nhận đã mua, số lượng thực tế, đơn giá & Người dùng & Trung bình\\ \hline
21 & UC21 & Chuyển mặt hàng đã mua vào tủ lạnh (*) & Tạo lô thực phẩm từ mặt hàng đã mua với hạn dùng đề xuất & Người dùng & Phức tạp\\ \hline
22 & UC22 & Sinh danh sách mua sắm từ kế hoạch bữa ăn (*) & Tính nguyên liệu còn thiếu cho các bữa đã lên lịch & Người dùng & Phức tạp\\ \hline
23 & UC23 & Xem lịch sử mua sắm & Xem các danh sách đã hoàn tất theo thời gian & Người dùng & Đơn giản\\ \hline
\multicolumn{6}{|l|}{\hs{Gói 4. Quản lý thực phẩm trong tủ lạnh}}\\ \hline
24 & UC24 & Xem danh sách thực phẩm trong tủ lạnh & Xem tồn kho theo vị trí, trạng thái hạn dùng & Người dùng & Đơn giản\\ \hline
25 & UC25 & Thêm thực phẩm vào tủ lạnh & Nhập lô thực phẩm với số lượng, hạn dùng, vị trí & Người dùng & Trung bình\\ \hline
26 & UC26 & Quét mã vạch sản phẩm (*) & Nhận diện thực phẩm đóng gói bằng camera & Người dùng, Cơ sở dữ liệu sản phẩm mở & Trung bình\\ \hline
27 & UC27 & Cập nhật thông tin thực phẩm & Sửa số lượng, hạn dùng, vị trí của lô & Người dùng & Đơn giản\\ \hline
28 & UC28 & Ghi nhận sử dụng thực phẩm & Trừ lượng đã dùng theo nguyên tắc FEFO & Người dùng & Trung bình\\ \hline
29 & UC29 & Loại bỏ thực phẩm & Ghi nhận thực phẩm hỏng hoặc hết hạn bị bỏ đi & Người dùng & Đơn giản\\ \hline
30 & UC30 & Xem thực phẩm sắp hết hạn & Xem các lô Sắp hết hạn, Đã hết hạn kèm món gợi ý & Người dùng & Đơn giản\\ \hline
31 & UC31 & Gửi nhắc nhở thực phẩm sắp hết hạn & Tác vụ hằng ngày gửi thông báo trước N ngày & Bộ lập lịch, Dịch vụ thông báo đẩy & Phức tạp\\ \hline
\multicolumn{6}{|l|}{\hs{Gói 5. Quản lý công thức nấu ăn}}\\ \hline
32 & UC32 & Thêm công thức nấu ăn & Tạo công thức với nguyên liệu có định lượng & Người dùng & Trung bình\\ \hline
33 & UC33 & Cập nhật, xóa công thức & Sửa hoặc xóa công thức trong bộ sưu tập & Người dùng & Đơn giản\\ \hline
34 & UC34 & Xem chi tiết công thức và đối chiếu nguyên liệu & Xem công thức, so sánh nguyên liệu với tủ lạnh & Người dùng & Trung bình\\ \hline
35 & UC35 & Đánh dấu công thức yêu thích & Thêm công thức vào hoặc bỏ khỏi danh sách yêu thích & Người dùng & Đơn giản\\ \hline
36 & UC36 & Bổ sung nguyên liệu còn thiếu vào danh sách mua sắm (*) & Đưa nguyên liệu thiếu của công thức vào danh sách & Người dùng & Trung bình\\ \hline
\multicolumn{6}{|l|}{\hs{Gói 6. Quản lý kế hoạch bữa ăn}}\\ \hline
37 & UC37 & Lên lịch bữa ăn & Thêm món vào bữa sáng, trưa, tối của một ngày & Người dùng & Trung bình\\ \hline
38 & UC38 & Xem lịch bữa ăn theo ngày, tuần & Xem kế hoạch, tình trạng nguyên liệu của từng món & Người dùng & Đơn giản\\ \hline
39 & UC39 & Cập nhật lịch bữa ăn & Đổi món, đổi khẩu phần, xóa món, di chuyển món & Người dùng & Đơn giản\\ \hline
40 & UC40 & Gợi ý món ăn từ thực phẩm sẵn có (*) & Xếp hạng món theo tồn kho, ưu tiên thực phẩm sắp hết hạn & Người dùng & Phức tạp\\ \hline
41 & UC41 & Xác nhận đã nấu món ăn (*) & Đánh dấu bữa đã nấu, trừ nguyên liệu theo FEFO & Người dùng & Phức tạp\\ \hline
\multicolumn{6}{|l|}{\hs{Gói 7. Tìm kiếm và phân loại thực phẩm}}\\ \hline
42 & UC42 & Tìm kiếm thực phẩm, món ăn & Tìm theo từ khóa, chấp nhận gõ không dấu & Người dùng & Trung bình\\ \hline
43 & UC43 & Lọc thực phẩm theo danh mục & Lọc theo loại, vị trí, trạng thái hạn dùng & Người dùng & Đơn giản\\ \hline
\multicolumn{6}{|l|}{\hs{Gói 8. Thống kê báo cáo}}\\ \hline
44 & UC44 & Xem báo cáo mua sắm và chi tiêu (*) & Thống kê chi tiêu, cơ cấu và mặt hàng mua nhiều & Người dùng & Trung bình\\ \hline
45 & UC45 & Xem báo cáo tiêu thụ thực phẩm & Thống kê lượng tiêu thụ theo thực phẩm, loại & Người dùng & Trung bình\\ \hline
46 & UC46 & Xem báo cáo lãng phí thực phẩm (*) & Thống kê giá trị, tỷ lệ thực phẩm bị bỏ & Người dùng & Trung bình\\ \hline
47 & UC47 & Xuất báo cáo (*) & Kết xuất báo cáo đang xem ra tệp PDF hoặc CSV & Người dùng & Đơn giản\\ \hline
\multicolumn{6}{|l|}{\hs{Gói 9. Quản trị hệ thống}}\\ \hline
48 & UC48 & Quản lý tài khoản người dùng & Tra cứu, khóa, mở khóa tài khoản & Quản trị viên & Trung bình\\ \hline
49 & UC49 & Quản lý danh mục thực phẩm & Thêm, sửa, ngừng sử dụng thực phẩm trong danh mục & Quản trị viên & Trung bình\\ \hline
50 & UC50 & Quản lý loại mặt hàng & Quản lý các loại thực phẩm dùng để phân loại & Quản trị viên & Đơn giản\\ \hline
51 & UC51 & Quản lý đơn vị đo lường & Quản lý đơn vị, nhóm đại lượng, hệ số quy đổi & Quản trị viên & Đơn giản\\ \hline
52 & UC52 & Quản lý công thức mẫu & Biên soạn công thức dùng chung cho mọi người dùng & Quản trị viên & Trung bình\\ \hline
53 & UC53 & Cấu hình tham số hệ thống & Điều chỉnh ngưỡng nhắc mặc định, giờ chạy, trọng số gợi ý & Quản trị viên & Trung bình\\ \hline
\end{longtable}
}

\subsection{Xác định các quan hệ}\label{ssec:quan-he-uc}

Mô hình sử dụng bốn loại quan hệ UML. Quan hệ \emph{kết hợp} (association) nối tác nhân với use case mà tác nhân tham gia; trên biểu đồ, tác nhân chính đặt bên trái ranh giới hệ thống và tác nhân phụ đặt bên phải. Quan hệ \emph{tổng quát hóa} giữa các tác nhân được trình bày tại mục~\ref{ssec:quan-he-tac-nhan}. Quan hệ \emph{bao gồm} (include) biểu thị hành vi luôn được thực hiện như một bước bắt buộc, chẳng hạn đăng ký tài khoản bao gồm xác thực OTP. Quan hệ \emph{mở rộng} (extend) biểu thị hành vi chỉ phát sinh trong một số điều kiện, chẳng hạn người dùng có thể chuyển mặt hàng vào tủ lạnh ngay sau khi đánh dấu đã mua. Bảng~\ref{tab:quan-he-uc} liệt kê các quan hệ include và extend trong mô hình.

{\small
\begin{longtable}{|L{3.6cm}|C{1.9cm}|L{3.4cm}|L{\dimexpr\textwidth-8\tabcolsep-5\arrayrulewidth-3.6cm-1.9cm-3.4cm\relax}|}
\caption{Các quan hệ bao gồm và mở rộng giữa các use case}\label{tab:quan-he-uc}\\
\hline
\hd{Use case nguồn} & \hd{Quan hệ} & \hd{Use case đích} & \hd{Ý nghĩa, điều kiện}\\ \hline
\endfirsthead
\multicolumn{4}{l}{\footnotesize\itshape Bảng \thetable\ (tiếp theo)}\\ \hline
\hd{Use case nguồn} & \hd{Quan hệ} & \hd{Use case đích} & \hd{Ý nghĩa, điều kiện}\\ \hline
\endhead
\uc{UC01} Đăng ký tài khoản & include & \uc{UC02} Xác thực mã OTP & Tài khoản chỉ được kích hoạt sau khi xác thực email\\ \hline
\uc{UC04} Khôi phục mật khẩu & include & \uc{UC02} Xác thực mã OTP & Mật khẩu chỉ được đặt lại sau khi xác thực email\\ \hline
\uc{UC13} Xóa thành viên & extend & \uc{UC12} Xem thông tin nhóm & Chỉ hiện với trưởng nhóm, tại dòng của từng thành viên\\ \hline
\uc{UC14} Chuyển quyền trưởng nhóm & extend & \uc{UC12} Xem thông tin nhóm & Chỉ hiện với trưởng nhóm khi nhóm có từ hai thành viên\\ \hline
\uc{UC14} Chuyển quyền trưởng nhóm & extend & \uc{UC15} Rời nhóm & Bắt buộc khi trưởng nhóm muốn rời nhóm còn thành viên khác (\br{BR07})\\ \hline
\uc{UC18} Sao chép danh sách & extend & \uc{UC16} Tạo danh sách & Người dùng chọn tạo danh sách từ một danh sách có sẵn\\ \hline
\uc{UC18} Sao chép danh sách & extend & \uc{UC23} Xem lịch sử mua sắm & Người dùng chọn “Mua lại” trên một danh sách đã hoàn tất\\ \hline
\uc{UC21} Chuyển vào tủ lạnh & extend & \uc{UC20} Đánh dấu đã mua & Người dùng đồng ý chuyển mặt hàng vừa mua vào tủ\\ \hline
\uc{UC22} Sinh danh sách từ kế hoạch & extend & \uc{UC38} Xem lịch bữa ăn & Người dùng chọn khoảng ngày và yêu cầu tạo danh sách\\ \hline
\uc{UC26} Quét mã vạch & extend & \uc{UC25} Thêm thực phẩm & Người dùng chọn nhập bằng camera thay vì gõ tên\\ \hline
\uc{UC27}, \uc{UC28}, \uc{UC29} & extend & \uc{UC24} Xem danh sách thực phẩm & Thao tác trên một lô được chọn trong tủ lạnh\\ \hline
\uc{UC33}, \uc{UC35}, \uc{UC36} & extend & \uc{UC34} Xem chi tiết công thức & Thao tác trên công thức đang xem; \uc{UC36} chỉ hiện khi có nguyên liệu thiếu\\ \hline
\uc{UC39}, \uc{UC41} & extend & \uc{UC38} Xem lịch bữa ăn & Thao tác trên một món trong lịch; \uc{UC41} chỉ áp dụng cho món chưa nấu\\ \hline
\uc{UC40} Gợi ý món ăn & extend & \uc{UC37} Lên lịch bữa ăn & Người dùng chọn món từ danh sách gợi ý thay vì tự tìm\\ \hline
\uc{UC43} Lọc theo danh mục & extend & \uc{UC42} Tìm kiếm & Người dùng áp bộ lọc lên kết quả tìm kiếm\\ \hline
\uc{UC34}, \uc{UC25} & extend & \uc{UC42} Tìm kiếm & Người dùng chọn một kết quả là món ăn hoặc thực phẩm\\ \hline
\uc{UC47} Xuất báo cáo & extend & \uc{UC44}, \uc{UC45}, \uc{UC46} & Người dùng chọn xuất báo cáo đang xem ra tệp\\ \hline
\uc{UC49} Quản lý danh mục thực phẩm & include & \uc{UC50}, \uc{UC51} & Mỗi thực phẩm phải thuộc một loại và có đơn vị mặc định\\ \hline
\end{longtable}
}

% ---------------------------------------------------------------------
%  2.2. Biểu đồ use case — 2.2.1 tổng quan, mở đầu 2.2.2
% ---------------------------------------------------------------------
\section{Biểu đồ use case}\label{sec:bieu-do-uc}

\subsection{Biểu đồ use case tổng quan}\label{ssec:uc-tong-quan}

Hình~\ref{fig:uc-tong-quan} thể hiện các chức năng của hệ thống và tác nhân tương ứng. Để biểu đồ dễ đọc, mỗi gói được biểu diễn bằng một use case tổng hợp, kèm mã các use case thành phần sẽ được phân rã tại mục~\ref{ssec:uc-phan-ra}. Quan hệ tổng quát hóa giữa các tác nhân đã trình bày ở Hình~\ref{fig:quan-he-tac-nhan} nên không lặp lại. Thành viên nhóm và Trưởng nhóm được đặt bên phải, gần các gói có liên kết riêng, để hạn chế giao cắt đường nối.

\diagram[0.82\textwidth]{c2-02-uc-tong-quan}{Biểu đồ use case tổng quan của hệ thống}{fig:uc-tong-quan}

Người dùng là tác nhân trung tâm, tham gia tám trong chín gói chức năng; ngoại lệ là gói Quản trị hệ thống. Thành viên nhóm và Trưởng nhóm kế thừa các liên kết của Người dùng, đồng thời có liên kết riêng với gói Quản lý nhóm gia đình. Trưởng nhóm còn tham gia gói Quản lý danh sách mua sắm để phân công công việc. Khách chỉ truy cập gói Xác thực và quản lý tài khoản, còn Quản trị viên sử dụng gói Quản trị hệ thống. Trong số các tác nhân phụ, Dịch vụ thư điện tử hỗ trợ các chức năng tài khoản; Dịch vụ thông báo đẩy gửi thông báo phân công, hoạt động nhóm và nhắc hạn dùng; Cơ sở dữ liệu sản phẩm mở phục vụ tra cứu mã vạch. Bộ lập lịch khởi chạy tác vụ nhắc hạn dùng hằng ngày mà không cần người dùng thao tác.

Ngoài các liên kết trên biểu đồ, các gói còn phụ thuộc nhau về dữ liệu. Gói danh sách mua sắm, tủ lạnh và kế hoạch bữa ăn cùng sử dụng dữ liệu của không gian hiện hành; gói công thức và tìm kiếm cung cấp dữ liệu tham chiếu cho các gói này. Gói thống kê chỉ đọc dữ liệu do các gói khác tạo ra, không làm thay đổi trạng thái hệ thống.

\subsection{Biểu đồ use case phân rã mức 2}\label{ssec:uc-phan-ra}

Mục này phân rã chín gói chức năng trong biểu đồ tổng quan. Mỗi gói gồm biểu đồ use case, phần giải thích các quan hệ và đặc tả các use case chính hoặc phức tạp theo mẫu của học phần. Những use case có nhiều nhánh xử lý hoặc thuật toán được bổ sung biểu đồ hoạt động; các use case phối hợp với hệ thống bên ngoài có thêm biểu đồ tuần tự. Với biểu mẫu nhập liệu, bảng dữ liệu đầu vào nêu rõ ràng buộc của từng trường.

% ---------------------------------------------------------------------
%  2.2.2.1. Phân rã gói "Xác thực và quản lý tài khoản"
% ---------------------------------------------------------------------
\subsubsection{Phân rã use case “Xác thực và quản lý tài khoản”}\label{sssec:pr-tai-khoan}

Gói Xác thực và quản lý tài khoản gồm tám use case từ \uc{UC01} đến \uc{UC08}, được phân rã tại Hình~\ref{fig:uc-tai-khoan}. Các use case \uc{UC01}, \uc{UC03} và \uc{UC04} dành cho Khách; \uc{UC05} đến \uc{UC08} dành cho Người dùng đã đăng nhập. \uc{UC02} Xác thực mã OTP được \uc{UC01} và \uc{UC04} bao gồm để dùng chung cơ chế xác minh email khi đăng ký hoặc khôi phục mật khẩu. Đây cũng là use case duy nhất kết nối trực tiếp với Dịch vụ thư điện tử trên biểu đồ. Use case mở rộng \uc{UC08} cho phép người dùng điều chỉnh ngưỡng nhắc hạn dùng và bật hoặc tắt từng loại thông báo.

\diagram[0.85\textwidth]{c2-03-uc-tai-khoan}{Biểu đồ use case phân rã gói Xác thực và quản lý tài khoản}{fig:uc-tai-khoan}

Trong \uc{UC01}, ứng dụng kiểm tra dữ liệu ngay trên thiết bị, kể cả khi ngoại tuyến, để phản hồi sớm và hạn chế yêu cầu không hợp lệ gửi lên máy chủ. Đặc tả, dữ liệu đầu vào và luồng hoạt động (bao gồm xác thực OTP) lần lượt được trình bày tại Bảng~\ref{tab:UC01}, Bảng~\ref{tab:in-dang-ky} và Hình~\ref{fig:act-dang-ky}.

\begin{ucspec}{UC01}{Đăng ký tài khoản}
\ucfield{Tác nhân}{Khách (chính); Dịch vụ thư điện tử (phụ, thông qua \uc{UC02})}
\ucfield{Mô tả}{Cho phép một người chưa có tài khoản tạo tài khoản mới bằng họ tên, email và mật khẩu, sau đó xác thực email để kích hoạt tài khoản.}
\ucfield{Điều kiện trước}{Khách chưa đăng nhập; ứng dụng đã được cài đặt trên thiết bị.}
\ucflow{Luồng thực thi chính}
\ucstep{1}{Khách}{Chọn chức năng “Đăng ký” trên màn hình chào.}
\ucstep{2}{Hệ thống}{Hiển thị biểu mẫu đăng ký gồm các trường trong Bảng~\ref{tab:in-dang-ky}.}
\ucstep{3}{Khách}{Nhập họ tên, email, mật khẩu, xác nhận mật khẩu, chọn đồng ý điều khoản sử dụng và nhấn “Đăng ký”.}
\ucstep{4}{Hệ thống}{Kiểm tra các trường bắt buộc và định dạng dữ liệu ngay trên thiết bị (\br{BR01}).}
\ucstep{5}{Hệ thống}{Gửi yêu cầu đăng ký lên máy chủ và kiểm tra email chưa được sử dụng.}
\ucstep{6}{Hệ thống}{Tạo tài khoản ở trạng thái Chờ xác thực, sinh mã OTP và yêu cầu Dịch vụ thư điện tử gửi mã tới email (\br{BR02}).}
\ucstep{7}{Hệ thống}{Chuyển sang màn hình nhập mã và thực hiện \uc{UC02} Xác thực mã OTP.}
\ucstep{8}{Hệ thống}{Khi \uc{UC02} thành công, kích hoạt tài khoản, tạo không gian cá nhân (\br{BR04}), tự động đăng nhập và hiển thị màn hình hướng dẫn nhanh.}
\ucflow{Luồng thực thi mở rộng}
\ucstep{4a}{Hệ thống}{Có trường bắt buộc bị bỏ trống: hiển thị \msg{MSG01} ngay dưới trường tương ứng; quay lại bước 3.}
\ucstep{4b}{Hệ thống}{Email sai định dạng hoặc mật khẩu không đủ mạnh: hiển thị \msg{MSG02} kèm yêu cầu cụ thể chưa đạt; quay lại bước 3.}
\ucstep{4c}{Hệ thống}{Mật khẩu xác nhận không khớp: hiển thị \msg{MSG29}; quay lại bước 3.}
\ucstep{5a}{Hệ thống}{Email đã thuộc một tài khoản đang hoạt động: hiển thị \msg{MSG03}, gợi ý đăng nhập hoặc khôi phục mật khẩu; kết thúc use case.}
\ucstep{5b}{Hệ thống}{Email thuộc một tài khoản còn Chờ xác thực: cập nhật thông tin đăng ký mới, sinh mã OTP mới và chuyển sang bước 7.}
\ucstep{5c}{Hệ thống}{Không kết nối được máy chủ: hiển thị \msg{MSG33}, giữ nguyên dữ liệu đã nhập để Khách thử lại.}
\ucstep{6a}{Hệ thống}{Dịch vụ thư điện tử không phản hồi: ghi nhật ký lỗi, vẫn chuyển sang bước 7 và cho phép yêu cầu gửi lại mã sau 60 giây.}
\ucfield{Điều kiện sau}{Thành công: tài khoản mới ở trạng thái Hoạt động, có một không gian cá nhân và người dùng đang ở trạng thái đăng nhập. Thất bại: không có tài khoản hoạt động nào được tạo.}
\ucfield{Quy tắc nghiệp vụ}{\br{BR01}, \br{BR02}, \br{BR04}}
\end{ucspec}

\begin{fieldtable}{Dữ liệu đầu vào của biểu mẫu đăng ký tài khoản}{tab:in-dang-ky}
\frow{Họ tên}{Tên hiển thị của người dùng trong nhóm}{Có}{Dài 2–50 ký tự, chỉ gồm chữ cái và khoảng trắng}{Nguyễn Văn An}
\frow{Email}{Định danh đăng nhập, nơi nhận mã OTP}{Có}{Đúng định dạng email, tối đa 100 ký tự, chưa thuộc tài khoản hoạt động}{an@gmail.com}
\frow{Mật khẩu}{Mật khẩu đăng nhập, hiển thị dạng ẩn}{Có}{Dài 8–32 ký tự, có chữ hoa, chữ thường và chữ số}{DiCho2025}
\frow{Xác nhận mật khẩu}{Nhập lại mật khẩu để tránh gõ nhầm}{Có}{Trùng khớp với trường Mật khẩu}{DiCho2025}
\frow{Đồng ý điều khoản}{Xác nhận đồng ý điều khoản sử dụng và chính sách quyền riêng tư}{Có}{Phải được chọn}{Đã chọn}
\end{fieldtable}

\diagram[0.92\textwidth]{c2-04-act-dang-ky}{Biểu đồ hoạt động của use case Đăng ký tài khoản và Xác thực mã OTP}{fig:act-dang-ky}

Hình~\ref{fig:act-dang-ky} chia luồng xử lý thành ba làn: Khách, ứng dụng trên thiết bị và máy chủ. Ứng dụng kiểm tra định dạng để phản hồi nhanh; máy chủ kiểm tra các điều kiện cần dữ liệu tập trung, như tính duy nhất của email và hiệu lực mã OTP. Giới hạn số lần nhập sai và thời hạn của mã giúp giảm nguy cơ dò mã tự động.

Trong \uc{UC03}, hệ thống xác thực thông tin đăng nhập và đăng ký mã thiết bị với Dịch vụ thông báo đẩy để gửi nhắc hạn dùng, thông báo phân công. Đặc tả của \uc{UC03} nằm ở Bảng~\ref{tab:UC03}; các use case \uc{UC02} và \uc{UC04} đến \uc{UC08} được đặc tả tại Phụ lục~\ref{pl:dac-ta-bo-sung}.

\begin{ucspec}{UC03}{Đăng nhập}
\ucfield{Tác nhân}{Khách}
\ucfield{Mô tả}{Cho phép người đã có tài khoản xác thực danh tính bằng email và mật khẩu để bắt đầu phiên làm việc.}
\ucfield{Điều kiện trước}{Khách đã có tài khoản; thiết bị có kết nối mạng.}
\ucflow{Luồng thực thi chính}
\ucstep{1}{Khách}{Mở ứng dụng và chọn “Đăng nhập”.}
\ucstep{2}{Hệ thống}{Hiển thị biểu mẫu gồm email, mật khẩu và tùy chọn “Ghi nhớ đăng nhập”.}
\ucstep{3}{Khách}{Nhập email, mật khẩu và nhấn “Đăng nhập”.}
\ucstep{4}{Hệ thống}{Kiểm tra các trường bắt buộc và định dạng email.}
\ucstep{5}{Hệ thống}{Đối chiếu thông tin đăng nhập với dữ liệu trên máy chủ.}
\ucstep{6}{Hệ thống}{Cấp mã truy cập và mã làm mới (\br{BR03}), đặt lại bộ đếm đăng nhập sai, đăng ký mã thiết bị với Dịch vụ thông báo đẩy.}
\ucstep{7}{Hệ thống}{Mở không gian làm việc gần nhất và hiển thị màn hình chính với ba thẻ tóm tắt: thực phẩm sắp hết hạn, danh sách mua sắm hôm nay và bữa ăn hôm nay.}
\ucflow{Luồng thực thi mở rộng}
\ucstep{3a}{Khách}{Chọn “Quên mật khẩu”: chuyển sang \uc{UC04} Khôi phục mật khẩu.}
\ucstep{4a}{Hệ thống}{Thiếu email hoặc mật khẩu: hiển thị \msg{MSG01}; quay lại bước 3.}
\ucstep{5a}{Hệ thống}{Sai email hoặc mật khẩu: hiển thị \msg{MSG05} mà không tiết lộ trường nào sai, tăng bộ đếm đăng nhập sai; quay lại bước 3.}
\ucstep{5b}{Hệ thống}{Sai lần thứ năm liên tiếp: khóa tạm tài khoản 15 phút và hiển thị \msg{MSG06} (\br{BR03}).}
\ucstep{5c}{Hệ thống}{Tài khoản bị quản trị viên khóa: hiển thị \msg{MSG21}; kết thúc use case.}
\ucstep{5d}{Hệ thống}{Tài khoản còn Chờ xác thực: gửi mã OTP mới và chuyển sang \uc{UC02}.}
\ucfield{Điều kiện sau}{Thành công: một phiên đăng nhập hợp lệ được tạo, Khách trở thành Người dùng. Nếu chọn “Ghi nhớ đăng nhập”, mã làm mới được lưu an toàn trên thiết bị để tự động đăng nhập lần sau.}
\ucfield{Quy tắc nghiệp vụ}{\br{BR03}}
\end{ucspec}

% ---------------------------------------------------------------------
%  2.2.2.2. Phân rã gói "Quản lý nhóm gia đình"
% ---------------------------------------------------------------------
\subsubsection{Phân rã use case “Quản lý nhóm gia đình”}\label{sssec:pr-nhom}

Gói Quản lý nhóm gia đình đáp ứng yêu cầu chia sẻ danh sách và phân vai trong gia đình. Hình~\ref{fig:uc-nhom} thể hiện các quyền theo vai trò: Người dùng chưa thuộc nhóm có thể tạo nhóm (\uc{UC09}) hoặc tham gia nhóm (\uc{UC11}); Thành viên nhóm có thể xem thông tin nhóm (\uc{UC12}) và rời nhóm (\uc{UC15}); Trưởng nhóm còn có quyền mời (\uc{UC10}), xóa thành viên (\uc{UC13}) và chuyển quyền (\uc{UC14}). Hai use case \uc{UC13} và \uc{UC14} mở rộng \uc{UC12} vì được thực hiện từ màn hình danh sách thành viên. Khi Trưởng nhóm rời nhóm còn thành viên khác, \uc{UC14} mở rộng \uc{UC15} để yêu cầu chọn người kế nhiệm, theo \br{BR05} và \br{BR07}.

\diagram[0.85\textwidth]{c2-05-uc-nhom}{Biểu đồ use case phân rã gói Quản lý nhóm gia đình}{fig:uc-nhom}

\begin{ucspec}{UC09}{Tạo nhóm gia đình}
\ucfield{Tác nhân}{Người dùng}
\ucfield{Mô tả}{Cho phép người dùng tạo một nhóm gia đình để chia sẻ danh sách mua sắm, tủ lạnh và kế hoạch bữa ăn với các thành viên khác; người tạo trở thành trưởng nhóm.}
\ucfield{Điều kiện trước}{Người dùng đã đăng nhập và chưa là thành viên của nhóm gia đình nào (\br{BR04}).}
\ucflow{Luồng thực thi chính}
\ucstep{1}{Người dùng}{Chọn “Tạo nhóm gia đình” tại màn hình Gia đình.}
\ucstep{2}{Hệ thống}{Hiển thị biểu mẫu gồm tên nhóm, ảnh nhóm (tùy chọn) và lựa chọn sao chép dữ liệu từ không gian cá nhân sang nhóm.}
\ucstep{3}{Người dùng}{Nhập tên nhóm, chọn ảnh, chọn có hoặc không sao chép dữ liệu, nhấn “Tạo nhóm”.}
\ucstep{4}{Hệ thống}{Kiểm tra tên nhóm dài 3–50 ký tự (\br{BR05}).}
\ucstep{5}{Hệ thống}{Tạo nhóm, ghi nhận người dùng là thành viên với vai trò trưởng nhóm.}
\ucstep{6}{Hệ thống}{Nếu người dùng chọn sao chép, sao chép các lô thực phẩm còn trong tủ và các danh sách mua sắm chưa hoàn tất từ không gian cá nhân sang nhóm.}
\ucstep{7}{Hệ thống}{Chuyển không gian làm việc sang nhóm mới, hiển thị \msg{MSG10} và gợi ý thực hiện \uc{UC10} để mời thành viên.}
\ucflow{Luồng thực thi mở rộng}
\ucstep{4a}{Hệ thống}{Tên nhóm không hợp lệ: hiển thị \msg{MSG02}; quay lại bước 3.}
\ucstep{5a}{Hệ thống}{Người dùng đã tham gia một nhóm gia đình khác trong lúc thao tác (ví dụ trên thiết bị khác): hiển thị \msg{MSG08}; kết thúc use case.}
\ucfield{Điều kiện sau}{Một nhóm gia đình mới tồn tại với đúng một thành viên là trưởng nhóm; không gian hiện hành của người dùng là nhóm này.}
\ucfield{Quy tắc nghiệp vụ}{\br{BR04}, \br{BR05}}
\end{ucspec}

Thay vì yêu cầu nhập email từng thành viên, hệ thống dùng mã mời có thời hạn. Trưởng nhóm có thể gửi mã qua tin nhắn hoặc để thành viên quét mã QR trực tiếp trên màn hình. Mã hết hạn sau 48 giờ và có thể bị thu hồi nếu bị lộ. Đặc tả use case mời và tham gia nhóm lần lượt nằm ở Bảng~\ref{tab:UC10} và Bảng~\ref{tab:UC11}.

\begin{ucspec}{UC10}{Mời thành viên vào nhóm}
\ucfield{Tác nhân}{Trưởng nhóm}
\ucfield{Mô tả}{Cho phép trưởng nhóm tạo mã mời gồm chữ và số cùng mã QR có thời hạn để chia sẻ cho người cần mời; trưởng nhóm cũng có thể thu hồi mã đang hiệu lực.}
\ucfield{Điều kiện trước}{Trưởng nhóm đã đăng nhập, không gian hiện hành là nhóm do mình làm trưởng nhóm.}
\ucflow{Luồng thực thi chính}
\ucstep{1}{Trưởng nhóm}{Chọn “Mời thành viên” tại màn hình thông tin nhóm.}
\ucstep{2}{Hệ thống}{Kiểm tra nhóm đã có mã mời còn hiệu lực hay chưa; nếu chưa, sinh mã mới gồm 8 ký tự theo \br{BR06} với thời hạn 48 giờ.}
\ucstep{3}{Hệ thống}{Hiển thị mã mời, mã QR tương ứng, thời điểm hết hạn và các nút “Sao chép”, “Chia sẻ”, “Tạo mã mới”, “Thu hồi”.}
\ucstep{4}{Trưởng nhóm}{Chọn “Chia sẻ”.}
\ucstep{5}{Hệ thống}{Mở bảng chia sẻ của hệ điều hành với nội dung lời mời chứa mã và liên kết mở ứng dụng.}
\ucflow{Luồng thực thi mở rộng}
\ucstep{1a}{Hệ thống}{Nhóm đã đủ số thành viên tối đa: hiển thị \msg{MSG09}; kết thúc use case.}
\ucstep{4a}{Trưởng nhóm}{Chọn “Tạo mã mới”: hệ thống thu hồi mã cũ, sinh mã mới và quay lại bước 3.}
\ucstep{4b}{Trưởng nhóm}{Chọn “Thu hồi”: hệ thống hiển thị \msg{MSG19} để xác nhận; khi được xác nhận, mã chuyển trạng thái Đã thu hồi và không thể dùng để tham gia nhóm.}
\ucfield{Điều kiện sau}{Nhóm có đúng một mã mời còn hiệu lực, hoặc không có mã nào nếu trưởng nhóm đã thu hồi.}
\ucfield{Quy tắc nghiệp vụ}{\br{BR05}, \br{BR06}}
\end{ucspec}

\begin{ucspec}{UC11}{Tham gia nhóm gia đình}
\ucfield{Tác nhân}{Người dùng}
\ucfield{Mô tả}{Cho phép người dùng gia nhập một nhóm gia đình bằng mã mời nhập tay hoặc bằng cách quét mã QR.}
\ucfield{Điều kiện trước}{Người dùng đã đăng nhập và đã nhận được mã mời từ trưởng nhóm.}
\ucflow{Luồng thực thi chính}
\ucstep{1}{Người dùng}{Chọn “Tham gia nhóm” tại màn hình Gia đình.}
\ucstep{2}{Hệ thống}{Hiển thị ô nhập mã mời và nút “Quét mã QR”.}
\ucstep{3}{Người dùng}{Nhập mã mời gồm 8 ký tự và nhấn “Tiếp tục”.}
\ucstep{4}{Hệ thống}{Tra cứu mã mời, kiểm tra mã tồn tại, còn hiệu lực và chưa bị thu hồi (\br{BR06}).}
\ucstep{5}{Hệ thống}{Kiểm tra người dùng chưa thuộc nhóm gia đình nào và nhóm chưa đủ số thành viên tối đa (\br{BR04}, \br{BR05}).}
\ucstep{6}{Hệ thống}{Hiển thị tên nhóm, tên trưởng nhóm, số thành viên hiện tại để người dùng xác nhận.}
\ucstep{7}{Người dùng}{Nhấn “Tham gia”.}
\ucstep{8}{Hệ thống}{Thêm người dùng vào nhóm với vai trò thành viên, chuyển không gian làm việc sang nhóm, gửi thông báo tới trưởng nhóm và hiển thị \msg{MSG10}.}
\ucflow{Luồng thực thi mở rộng}
\ucstep{3a}{Người dùng}{Chọn “Quét mã QR”: hệ thống xin quyền truy cập camera, đọc mã từ QR và chuyển sang bước 4.}
\ucstep{3b}{Người dùng}{Mở liên kết mời từ tin nhắn: hệ điều hành mở ứng dụng với mã đã điền sẵn và chuyển sang bước 4.}
\ucstep{4a}{Hệ thống}{Mã không tồn tại, hết hạn hoặc đã bị thu hồi: hiển thị \msg{MSG07}; quay lại bước 3.}
\ucstep{5a}{Hệ thống}{Người dùng đang thuộc một nhóm gia đình khác: hiển thị \msg{MSG08}; kết thúc use case.}
\ucstep{5b}{Hệ thống}{Nhóm đã đủ số thành viên: hiển thị \msg{MSG09}; kết thúc use case.}
\ucstep{7a}{Người dùng}{Nhấn “Hủy”: kết thúc use case, không có thay đổi.}
\ucfield{Điều kiện sau}{Người dùng là thành viên của nhóm, có quyền đọc và ghi dữ liệu chung của nhóm; dữ liệu trong không gian cá nhân được giữ nguyên.}
\ucfield{Quy tắc nghiệp vụ}{\br{BR04}, \br{BR05}, \br{BR06}, \br{BR23}}
\end{ucspec}

\diagram[0.7\textwidth]{c2-06-act-tham-gia-nhom}{Biểu đồ hoạt động của use case Tham gia nhóm gia đình}{fig:act-tham-gia-nhom}

Biểu đồ hoạt động ở Hình~\ref{fig:act-tham-gia-nhom} thể hiện thứ tự kiểm tra ba điều kiện: mã còn hiệu lực, người dùng chưa là thành viên của nhóm khác và nhóm còn chỗ. Hệ thống kiểm tra các điều kiện này trước khi hiển thị bước xác nhận, tránh yêu cầu người dùng tiếp tục khi thao tác sẽ bị từ chối. Các use case \uc{UC12} đến \uc{UC15} được đặc tả tại Phụ lục~\ref{pl:dac-ta-bo-sung}.

% ---------------------------------------------------------------------
%  2.2.2.3. Phân rã gói "Quản lý danh sách mua sắm"
% ---------------------------------------------------------------------
\subsubsection{Phân rã use case “Quản lý danh sách mua sắm”}\label{sssec:pr-mua-sam}

Gói Quản lý danh sách mua sắm gồm tám use case, được phân rã tại Hình~\ref{fig:uc-mua-sam}. Người dùng có thể tạo danh sách, cập nhật mặt hàng, đánh dấu đã mua, sinh danh sách từ kế hoạch bữa ăn và xem lịch sử. Riêng \uc{UC19} Phân công người mua chỉ dành cho Trưởng nhóm và gửi thông báo qua Dịch vụ thông báo đẩy. Use case \uc{UC18} mở rộng cả \uc{UC16} và \uc{UC23}, cho phép sao chép danh sách từ hai điểm vào. Use case \uc{UC21} mở rộng \uc{UC20}, liên kết trực tiếp gói mua sắm với gói tủ lạnh.

\diagram[0.85\textwidth]{c2-07-uc-mua-sam}{Biểu đồ use case phân rã gói Quản lý danh sách mua sắm}{fig:uc-mua-sam}

\begin{ucspec}{UC16}{Tạo danh sách mua sắm}
\ucfield{Tác nhân}{Người dùng}
\ucfield{Mô tả}{Cho phép người dùng tạo một danh sách mua sắm gắn với ngày dự kiến đi chợ trong không gian hiện hành.}
\ucfield{Điều kiện trước}{Người dùng đã đăng nhập.}
\ucflow{Luồng thực thi chính}
\ucstep{1}{Người dùng}{Tại màn hình Mua sắm, nhấn “Tạo danh sách”.}
\ucstep{2}{Hệ thống}{Hiển thị biểu mẫu với tên danh sách gợi ý theo ngày (ví dụ “Đi chợ thứ Bảy 13/12”), ngày dự kiến mua mặc định là hôm nay và ô ghi chú.}
\ucstep{3}{Người dùng}{Điều chỉnh tên, ngày dự kiến mua, ghi chú và nhấn “Tạo”.}
\ucstep{4}{Hệ thống}{Kiểm tra tên dài 1–100 ký tự và ngày dự kiến mua không ở quá khứ (\br{BR08}).}
\ucstep{5}{Hệ thống}{Tạo danh sách ở trạng thái Đang lập trong không gian hiện hành, đồng bộ tới các thành viên khác của nhóm.}
\ucstep{6}{Hệ thống}{Mở màn hình chi tiết danh sách để người dùng thêm mặt hàng (\uc{UC17}).}
\ucflow{Luồng thực thi mở rộng}
\ucstep{1a}{Người dùng}{Chọn “Tạo từ danh sách cũ”: thực hiện \uc{UC18}.}
\ucstep{4a}{Hệ thống}{Ngày dự kiến mua ở quá khứ: hiển thị \msg{MSG02}; quay lại bước 3.}
\ucstep{5a}{Hệ thống}{Mất kết nối mạng: lưu danh sách trên thiết bị, hiển thị \msg{MSG20} và đồng bộ khi có mạng trở lại (\nfr{NFR11}).}
\ucfield{Điều kiện sau}{Một danh sách mua sắm mới ở trạng thái Đang lập thuộc không gian hiện hành.}
\ucfield{Quy tắc nghiệp vụ}{\br{BR04}, \br{BR08}}
\end{ucspec}

\uc{UC17} xử lý việc gộp mặt hàng trùng. Khi người dùng thêm thực phẩm đã có trong danh sách, hệ thống quy đổi số lượng mới về đơn vị của mặt hàng hiện tại rồi cộng dồn. Chẳng hạn, “Thịt ba chỉ 3 lạng” và “Thịt ba chỉ 0,5 kg” được gộp thành một dòng “Thịt ba chỉ 8 lạng”. Đặc tả use case nằm ở Bảng~\ref{tab:UC17}; các trường dữ liệu của mặt hàng được liệt kê tại Bảng~\ref{tab:in-mat-hang}.

\begin{ucspec}{UC17}{Cập nhật mặt hàng trong danh sách}
\ucfield{Tác nhân}{Người dùng}
\ucfield{Mô tả}{Cho phép người dùng thêm, sửa hoặc xóa mặt hàng trong một danh sách mua sắm chưa hoàn tất; hệ thống tự động gộp các mặt hàng trùng.}
\ucfield{Điều kiện trước}{Người dùng đang xem một danh sách mua sắm ở trạng thái Đang lập hoặc Đang mua.}
\ucflow{Luồng thực thi chính}
\ucstep{1}{Người dùng}{Gõ tên thực phẩm vào ô “Thêm mặt hàng”.}
\ucstep{2}{Hệ thống}{Gợi ý tức thì các thực phẩm trong danh mục khớp với từ khóa, chấp nhận gõ không dấu, kèm đơn vị mặc định.}
\ucstep{3}{Người dùng}{Chọn thực phẩm, nhập số lượng, chọn đơn vị, ghi chú (nếu có) và nhấn “Thêm”.}
\ucstep{4}{Hệ thống}{Kiểm tra số lượng là số dương với tối đa hai chữ số thập phân (\br{BR14}).}
\ucstep{5}{Hệ thống}{Kiểm tra danh sách đã có mặt hàng cùng thực phẩm và cùng nhóm đại lượng đơn vị hay chưa (\br{BR09}).}
\ucstep{6}{Hệ thống}{Nếu chưa có, tạo mặt hàng mới ở trạng thái Cần mua; hiển thị mặt hàng trong danh sách, nhóm theo loại thực phẩm.}
\ucstep{7}{Hệ thống}{Đồng bộ thay đổi tới các thành viên khác đang mở danh sách.}
\ucflow{Luồng thực thi mở rộng}
\ucstep{1a}{Người dùng}{Vuốt sang trái trên một mặt hàng để sửa hoặc xóa; với thao tác xóa, hệ thống cho phép hoàn tác trong 5 giây thay vì hỏi xác nhận.}
\ucstep{1b}{Hệ thống}{Mặt hàng đã ở trạng thái Đã nhập tủ lạnh: không cho phép sửa hoặc xóa, hiển thị \msg{MSG32}.}
\ucstep{3a}{Người dùng}{Thực phẩm không có trong danh mục: chọn “Thêm ‘<tên đã gõ>’” để tạo mặt hàng với tên tự do; mặt hàng này vẫn được lưu nhưng không tham gia đối chiếu tồn kho.}
\ucstep{4a}{Hệ thống}{Số lượng không hợp lệ: hiển thị \msg{MSG16}; quay lại bước 3.}
\ucstep{6a}{Hệ thống}{Đã có mặt hàng trùng: quy đổi và cộng dồn số lượng vào mặt hàng hiện có (\br{BR16}), hiển thị \msg{MSG11}.}
\ucstep{6b}{Hệ thống}{Mặt hàng trùng nhưng khác nhóm đại lượng đơn vị (ví dụ “2 bó” và “300 g”): tạo dòng riêng và đánh dấu biểu tượng cảnh báo trùng tên.}
\ucfield{Điều kiện sau}{Danh sách phản ánh đúng các mặt hàng cần mua, không có hai dòng trùng thực phẩm cùng nhóm đại lượng; mọi thành viên thấy cùng một trạng thái.}
\ucfield{Quy tắc nghiệp vụ}{\br{BR09}, \br{BR14}, \br{BR16}}
\end{ucspec}

\begin{fieldtable}{Dữ liệu đầu vào của một mặt hàng mua sắm}{tab:in-mat-hang}
\frow{Thực phẩm}{Tham chiếu tới danh mục thực phẩm, hoặc tên tự do}{Có}{Chọn từ gợi ý hoặc nhập tên dài 1–100 ký tự}{Cà chua}
\frow{Số lượng}{Lượng cần mua}{Có}{Số dương, tối đa hai chữ số thập phân}{1,5}
\frow{Đơn vị}{Đơn vị đo của số lượng}{Có}{Thuộc danh mục đơn vị; mặc định là đơn vị của thực phẩm}{kg}
\frow{Ghi chú}{Yêu cầu cụ thể khi mua}{Không}{Tối đa 200 ký tự}{Chọn quả chín vừa}
\frow{Người được giao}{Thành viên phụ trách mua}{Không}{Do trưởng nhóm chỉ định (\uc{UC19})}{Đinh Đức}
\frow{Đơn giá}{Giá thực tế khi mua, nhập ở \uc{UC20}}{Không}{Số không âm, đơn vị đồng}{25.000}
\end{fieldtable}

Phân công người mua là một trong ba use case phức tạp của gói vì liên quan đến quyền hạn, trạng thái mặt hàng và dịch vụ bên ngoài. Trưởng nhóm có thể giao từng mặt hàng hoặc chọn nhiều mặt hàng cho cùng một thành viên, chẳng hạn giao rau củ cho người đi chợ gần nhà và đồ khô cho người đi siêu thị.

\begin{ucspec}{UC19}{Phân công người mua}
\ucfield{Tác nhân}{Trưởng nhóm (chính); Dịch vụ thông báo đẩy (phụ)}
\ucfield{Mô tả}{Cho phép trưởng nhóm giao một hoặc nhiều mặt hàng cho một thành viên phụ trách mua; người được giao nhận thông báo đẩy.}
\ucfield{Điều kiện trước}{Trưởng nhóm đang xem một danh sách chưa hoàn tất của nhóm; người được giao có thể là một thành viên khác hoặc chính trưởng nhóm.}
\ucflow{Luồng thực thi chính}
\ucstep{1}{Trưởng nhóm}{Nhấn giữ một mặt hàng để vào chế độ chọn nhiều, chọn thêm các mặt hàng khác, nhấn “Phân công”.}
\ucstep{2}{Hệ thống}{Hiển thị danh sách thành viên kèm số mặt hàng mỗi người đang được giao trong danh sách này.}
\ucstep{3}{Trưởng nhóm}{Chọn một thành viên và nhấn “Giao việc”.}
\ucstep{4}{Hệ thống}{Kiểm tra quyền trưởng nhóm và kiểm tra các mặt hàng được chọn đều ở trạng thái Cần mua hoặc Đã phân công (\br{BR10}).}
\ucstep{5}{Hệ thống}{Ghi nhận người được giao, chuyển các mặt hàng sang Đã phân công, đồng bộ tới các thành viên.}
\ucstep{6}{Hệ thống}{Tạo thông báo “Bạn được giao mua k mặt hàng trong danh sách <tên>” và yêu cầu Dịch vụ thông báo đẩy gửi tới thiết bị của người được giao.}
\ucstep{7}{Hệ thống}{Hiển thị \msg{MSG10}; mỗi mặt hàng hiển thị ảnh đại diện của người được giao.}
\ucflow{Luồng thực thi mở rộng}
\ucstep{3a}{Trưởng nhóm}{Chọn “Bỏ phân công”: hệ thống chuyển các mặt hàng về Cần mua và thông báo cho người được giao trước đó.}
\ucstep{4a}{Hệ thống}{Người thực hiện không còn là trưởng nhóm (do vừa chuyển quyền trên thiết bị khác): hiển thị \msg{MSG13}; kết thúc use case.}
\ucstep{4b}{Hệ thống}{Có mặt hàng đã ở trạng thái Đã mua hoặc Đã bỏ qua: bỏ qua các mặt hàng đó, chỉ phân công các mặt hàng hợp lệ và thông báo số mặt hàng bị bỏ qua.}
\ucstep{6a}{Hệ thống}{Gửi thông báo đẩy thất bại: thông báo vẫn được lưu trong hộp thư trong ứng dụng, hệ thống thử gửi lại tối đa 3 lần (\nfr{NFR12}).}
\ucfield{Điều kiện sau}{Các mặt hàng được chọn có đúng một người được giao; người được giao nhận được thông báo trong hộp thư và trên thiết bị.}
\ucfield{Quy tắc nghiệp vụ}{\br{BR07}, \br{BR10}}
\end{ucspec}

\uc{UC20} và \uc{UC21} thường được sử dụng ngay tại điểm mua sắm. Thao tác chính được rút gọn thành một lần chạm vào ô đánh dấu; các thông tin bổ sung như số lượng thực tế, đơn giá và việc chuyển thực phẩm vào tủ lạnh là tùy chọn, với giá trị mặc định được điền sẵn. Luồng phối hợp giữa hai use case được thể hiện ở Hình~\ref{fig:act-danh-dau-da-mua}.

\begin{ucspec}{UC20}{Đánh dấu mặt hàng đã mua}
\ucfield{Tác nhân}{Người dùng}
\ucfield{Mô tả}{Cho phép người dùng ghi nhận một mặt hàng đã được mua, kèm số lượng thực tế và đơn giá, để các thành viên khác không mua trùng.}
\ucfield{Điều kiện trước}{Danh sách mua sắm chưa hoàn tất; mặt hàng ở trạng thái Cần mua hoặc Đã phân công.}
\ucflow{Luồng thực thi chính}
\ucstep{1}{Người dùng}{Chạm vào ô đánh dấu bên cạnh mặt hàng.}
\ucstep{2}{Hệ thống}{Kiểm tra mặt hàng có đang được giao cho người khác hay không (\br{BR10}).}
\ucstep{3}{Hệ thống}{Chuyển mặt hàng sang Đã mua, ghi người mua và thời điểm mua, gạch ngang mặt hàng và đưa xuống cuối danh sách.}
\ucstep{4}{Hệ thống}{Hiển thị thanh nhập nhanh cho phép sửa số lượng thực tế và nhập đơn giá, cùng lựa chọn “Cho vào tủ lạnh”.}
\ucstep{5}{Người dùng}{Nhập đơn giá (nếu muốn) và đóng thanh nhập nhanh.}
\ucstep{6}{Hệ thống}{Lưu dữ liệu, cập nhật trạng thái danh sách sang Đang mua hoặc Hoàn tất theo \br{BR08}, đồng bộ tới các thành viên.}
\ucflow{Luồng thực thi mở rộng}
\ucstep{2a}{Hệ thống}{Mặt hàng đang được giao cho thành viên khác: hiển thị \msg{MSG12}; nếu người dùng không đồng ý thì kết thúc use case, nếu đồng ý thì tiếp tục bước 3 và gửi thông báo cho người được giao.}
\ucstep{5a}{Người dùng}{Chọn “Cho vào tủ lạnh”: thực hiện \uc{UC21}.}
\ucstep{5b}{Người dùng}{Chạm lại vào ô đánh dấu trước khi mặt hàng được nhập tủ: hoàn tác, mặt hàng trở về trạng thái Cần mua.}
\ucstep{6a}{Hệ thống}{Mất kết nối mạng: lưu thay đổi trên thiết bị, hiển thị biểu tượng chờ đồng bộ; khi có mạng, nếu mặt hàng đã được người khác đánh dấu trước thì giữ bản ghi sớm hơn và thông báo cho người dùng (\nfr{NFR11}).}
\ucfield{Điều kiện sau}{Mặt hàng ở trạng thái Đã mua với người mua, thời điểm, số lượng thực tế và đơn giá (nếu có); trạng thái danh sách được cập nhật.}
\ucfield{Quy tắc nghiệp vụ}{\br{BR08}, \br{BR10}, \br{BR14}}
\end{ucspec}

\begin{ucspec}{UC21}{Chuyển mặt hàng đã mua vào tủ lạnh}
\ucfield{Tác nhân}{Người dùng}
\ucfield{Mô tả}{Cho phép người dùng chuyển một hoặc nhiều mặt hàng đã mua thành các lô thực phẩm trong tủ lạnh mà không phải nhập lại thông tin.}
\ucfield{Điều kiện trước}{Có ít nhất một mặt hàng ở trạng thái Đã mua tham chiếu tới thực phẩm trong danh mục.}
\ucflow{Luồng thực thi chính}
\ucstep{1}{Người dùng}{Chọn “Cho vào tủ lạnh” trên một mặt hàng, hoặc chọn “Cất tất cả vào tủ” sau khi đi chợ về.}
\ucstep{2}{Hệ thống}{Với mỗi mặt hàng, điền sẵn tên, số lượng thực tế, đơn vị, ngày mua bằng ngày đánh dấu đã mua, vị trí lưu trữ mặc định theo loại thực phẩm và ngày hết hạn đề xuất bằng ngày mua cộng hạn dùng gợi ý tại vị trí đó (\br{BR11}).}
\ucstep{3}{Người dùng}{Rà soát, điều chỉnh vị trí hoặc ngày hết hạn theo bao bì (nếu cần) và nhấn “Xác nhận”.}
\ucstep{4}{Hệ thống}{Kiểm tra ngày hết hạn không trước ngày mua (\br{BR11}).}
\ucstep{5}{Hệ thống}{Tạo các lô thực phẩm trong tủ của không gian hiện hành, liên kết mỗi lô với mặt hàng nguồn để lưu đơn giá; tính trạng thái hạn dùng ban đầu (\br{BR12}).}
\ucstep{6}{Hệ thống}{Chuyển các mặt hàng sang Đã nhập tủ lạnh, hiển thị \msg{MSG10} kèm lối tắt tới màn hình Tủ lạnh.}
\ucflow{Luồng thực thi mở rộng}
\ucstep{2a}{Hệ thống}{Mặt hàng có tên tự do không thuộc danh mục: yêu cầu người dùng chọn thực phẩm tương ứng hoặc bỏ qua mặt hàng đó.}
\ucstep{3a}{Người dùng}{Thay đổi vị trí lưu trữ: hệ thống tính lại ngày hết hạn đề xuất theo vị trí mới, trừ khi người dùng đã sửa tay ngày hết hạn.}
\ucstep{4a}{Hệ thống}{Ngày hết hạn trước ngày mua: hiển thị \msg{MSG14} tại dòng tương ứng; quay lại bước 3.}
\ucstep{5a}{Hệ thống}{Ngày hết hạn đã nằm trong ngưỡng nhắc: lô được tạo với trạng thái Sắp hết hạn và được hiển thị nổi bật trên màn hình kết quả.}
\ucfield{Điều kiện sau}{Mỗi mặt hàng được chọn sinh ra đúng một lô thực phẩm trong tủ; mặt hàng ở trạng thái Đã nhập tủ lạnh và không thể hoàn tác.}
\ucfield{Quy tắc nghiệp vụ}{\br{BR11}, \br{BR12}}
\end{ucspec}

\diagram[0.66\textwidth]{c2-08-act-danh-dau-da-mua}{Biểu đồ hoạt động của use case Đánh dấu đã mua và Chuyển vào tủ lạnh}{fig:act-danh-dau-da-mua}

\uc{UC22} liên kết kế hoạch bữa ăn với danh sách mua sắm. Thuật toán ở Hình~\ref{fig:act-sinh-ds} tính lượng cần mua bằng cách lấy tổng nhu cầu trừ lượng tồn kho còn hạn đến ngày nấu và lượng đã có trong các danh sách chưa hoàn tất. Cách tính này tránh đưa trùng nguyên liệu vào danh sách khi các khoảng ngày được lập kế hoạch chồng lấn.

\begin{ucspec}{UC22}{Sinh danh sách mua sắm từ kế hoạch bữa ăn}
\ucfield{Tác nhân}{Người dùng}
\ucfield{Mô tả}{Cho phép người dùng tạo danh sách mua sắm chứa đúng phần nguyên liệu còn thiếu cho các bữa ăn đã lên lịch trong một khoảng ngày.}
\ucfield{Điều kiện trước}{Lịch bữa ăn trong khoảng ngày được chọn có ít nhất một món chưa nấu gắn với công thức.}
\ucflow{Luồng thực thi chính}
\ucstep{1}{Người dùng}{Tại màn hình lịch bữa ăn, chọn khoảng ngày (mặc định là tuần đang xem) và nhấn “Tạo danh sách mua sắm”.}
\ucstep{2}{Hệ thống}{Lấy các món chưa nấu trong khoảng ngày, nhân định lượng nguyên liệu theo tỷ lệ số khẩu phần dự kiến trên số khẩu phần chuẩn của công thức.}
\ucstep{3}{Hệ thống}{Gộp nhu cầu theo thực phẩm sau khi quy đổi về đơn vị cơ sở của cùng nhóm đại lượng (\br{BR16}).}
\ucstep{4}{Hệ thống}{Với mỗi nguyên liệu, tính lượng thiếu theo \br{BR20}.}
\ucstep{5}{Hệ thống}{Hiển thị danh sách đề xuất gồm các nguyên liệu có lượng thiếu dương, mỗi dòng ghi rõ món ăn sử dụng và ngày cần dùng sớm nhất; ngày dự kiến mua đề xuất là ngày liền trước bữa sớm nhất nhưng không sớm hơn hôm nay (\br{BR08}).}
\ucstep{6}{Người dùng}{Bỏ chọn các nguyên liệu không muốn mua, chỉnh số lượng và nhấn “Tạo danh sách”.}
\ucstep{7}{Hệ thống}{Tạo danh sách mới ở trạng thái Đang lập với các mặt hàng đã chọn, hiển thị \msg{MSG10}.}
\ucflow{Luồng thực thi mở rộng}
\ucstep{2a}{Hệ thống}{Có món không gắn công thức (món nhập tên tự do): bỏ qua món đó và ghi chú trong danh sách đề xuất.}
\ucstep{5a}{Hệ thống}{Không có nguyên liệu nào thiếu: hiển thị \msg{MSG15}; kết thúc use case.}
\ucstep{6a}{Người dùng}{Chọn “Gộp vào danh sách có sẵn”: hệ thống liệt kê các danh sách chưa hoàn tất có ngày dự kiến mua trong khoảng, người dùng chọn một danh sách; hệ thống gộp mặt hàng theo \br{BR09}.}
\ucfield{Điều kiện sau}{Một danh sách mua sắm chứa phần nguyên liệu còn thiếu được tạo mới hoặc được gộp vào danh sách có sẵn.}
\ucfield{Quy tắc nghiệp vụ}{\br{BR09}, \br{BR16}, \br{BR20}}
\end{ucspec}

\diagram[0.64\textwidth]{c2-09-act-sinh-ds-tu-ke-hoach}{Biểu đồ hoạt động của use case Sinh danh sách mua sắm từ kế hoạch bữa ăn}{fig:act-sinh-ds}

Các use case \uc{UC18} Sao chép danh sách mua sắm đã có và \uc{UC23} Xem lịch sử mua sắm được đặc tả tại Phụ lục~\ref{pl:dac-ta-bo-sung}.

% ---------------------------------------------------------------------
%  2.2.2.4. Phân rã gói "Quản lý thực phẩm trong tủ lạnh"
% ---------------------------------------------------------------------
\subsubsection{Phân rã use case “Quản lý thực phẩm trong tủ lạnh”}\label{sssec:pr-tu-lanh}

Gói Quản lý thực phẩm trong tủ lạnh gồm tám use case từ \uc{UC24} đến \uc{UC31}, được phân rã tại Hình~\ref{fig:uc-tu-lanh}. \uc{UC24} Xem danh sách thực phẩm là điểm vào chính. Các thao tác cập nhật (\uc{UC27}), ghi nhận sử dụng (\uc{UC28}) và loại bỏ (\uc{UC29}) mở rộng use case này, vì được thực hiện trên từng lô trong danh sách. \uc{UC26} Quét mã vạch mở rộng \uc{UC25} như một phương thức nhập liệu và là use case duy nhất tra cứu Cơ sở dữ liệu sản phẩm mở. Riêng \uc{UC31} được Bộ lập lịch khởi phát thay vì người dùng.

\diagram[0.8\textwidth]{c2-10-uc-tu-lanh}{Biểu đồ use case phân rã gói Quản lý thực phẩm trong tủ lạnh}{fig:uc-tu-lanh}

\begin{ucspec}{UC25}{Thêm thực phẩm vào tủ lạnh}
\ucfield{Tác nhân}{Người dùng}
\ucfield{Mô tả}{Cho phép người dùng nhập một lô thực phẩm vào tủ lạnh của không gian hiện hành với số lượng, ngày mua, ngày hết hạn và vị trí lưu trữ.}
\ucfield{Điều kiện trước}{Người dùng đã đăng nhập.}
\ucflow{Luồng thực thi chính}
\ucstep{1}{Người dùng}{Tại màn hình Tủ lạnh, nhấn nút “Thêm thực phẩm”.}
\ucstep{2}{Hệ thống}{Hiển thị biểu mẫu gồm các trường trong Bảng~\ref{tab:in-thuc-pham} và nút “Quét mã vạch”.}
\ucstep{3}{Người dùng}{Gõ tên và chọn thực phẩm từ danh sách gợi ý.}
\ucstep{4}{Hệ thống}{Điền sẵn đơn vị mặc định, vị trí lưu trữ mặc định theo loại thực phẩm, ngày mua là hôm nay và ngày hết hạn đề xuất theo \br{BR11}.}
\ucstep{5}{Người dùng}{Nhập số lượng, điều chỉnh các trường còn lại nếu cần và nhấn “Lưu”.}
\ucstep{6}{Hệ thống}{Kiểm tra các trường bắt buộc, số lượng dương (\br{BR14}) và ngày hết hạn không trước ngày mua (\br{BR11}).}
\ucstep{7}{Hệ thống}{Tạo lô thực phẩm, tính trạng thái hạn dùng theo \br{BR12}, hiển thị lô trong nhóm vị trí tương ứng và thông báo \msg{MSG10}.}
\ucflow{Luồng thực thi mở rộng}
\ucstep{2a}{Người dùng}{Nhấn “Quét mã vạch”: thực hiện \uc{UC26}; khi nhận diện thành công, hệ thống điền sẵn thực phẩm và chuyển sang bước 4.}
\ucstep{3a}{Người dùng}{Thực phẩm không có trong danh mục: chọn “Thêm thực phẩm mới”, nhập tên và chọn loại; thực phẩm được lưu dưới dạng thực phẩm riêng của không gian, chờ quản trị viên xem xét bổ sung vào danh mục chung.}
\ucstep{6a}{Hệ thống}{Thiếu trường bắt buộc: hiển thị \msg{MSG01}; quay lại bước 5.}
\ucstep{6b}{Hệ thống}{Số lượng không hợp lệ: hiển thị \msg{MSG16}; quay lại bước 5.}
\ucstep{6c}{Hệ thống}{Ngày hết hạn trước ngày mua: hiển thị \msg{MSG14}; quay lại bước 5.}
\ucstep{7a}{Hệ thống}{Ngày hết hạn đã qua: vẫn lưu lô với trạng thái Đã hết hạn sau khi người dùng xác nhận, nhằm phản ánh đúng thực tế trong tủ.}
\ucfield{Điều kiện sau}{Một lô thực phẩm mới thuộc không gian hiện hành, có trạng thái hạn dùng được tính đúng; lô được đưa vào phạm vi kiểm tra của \uc{UC31}.}
\ucfield{Quy tắc nghiệp vụ}{\br{BR11}, \br{BR12}, \br{BR14}, \br{BR16}}
\end{ucspec}

\begin{fieldtable}{Dữ liệu đầu vào của một lô thực phẩm trong tủ lạnh}{tab:in-thuc-pham}
\frow{Thực phẩm}{Tham chiếu tới danh mục thực phẩm}{Có}{Chọn từ gợi ý, từ kết quả quét mã vạch hoặc tạo thực phẩm riêng}{Sữa tươi không đường}
\frow{Số lượng}{Lượng thực phẩm cất vào tủ}{Có}{Số dương, tối đa hai chữ số thập phân}{2}
\frow{Đơn vị}{Đơn vị đo của số lượng}{Có}{Thuộc danh mục đơn vị; mặc định theo thực phẩm}{hộp}
\frow{Vị trí lưu trữ}{Nơi cất trong tủ hoặc trong bếp}{Có}{Một trong năm vị trí của Bảng~\ref{tab:vi-tri-luu-tru}}{Cánh tủ}
\frow{Ngày mua}{Ngày đưa thực phẩm vào tủ}{Có}{Định dạng dd/MM/yyyy, không sau ngày hiện tại}{12/12/2025}
\frow{Ngày hết hạn}{Ngày thực phẩm hết hạn sử dụng}{Có}{Không trước ngày mua; đề xuất tự động nếu bỏ trống}{19/12/2025}
\frow{Ghi chú}{Thông tin bổ sung}{Không}{Tối đa 200 ký tự}{Đã mở nắp 1 hộp}
\end{fieldtable}

\uc{UC28} ghi nhận lượng thực phẩm đã dùng nhưng không gắn với bữa ăn trong kế hoạch. Người dùng nhập lượng ở cấp thực phẩm; hệ thống tự phân bổ phần giảm trừ cho các lô theo nguyên tắc FEFO. Nhờ đó, người dùng không cần chọn lô cụ thể.

\begin{ucspec}{UC28}{Ghi nhận sử dụng thực phẩm}
\ucfield{Tác nhân}{Người dùng}
\ucfield{Mô tả}{Cho phép người dùng ghi nhận lượng thực phẩm đã lấy ra sử dụng; hệ thống trừ tồn kho theo nguyên tắc hết hạn trước dùng trước.}
\ucfield{Điều kiện trước}{Thực phẩm có ít nhất một lô với lượng còn lại lớn hơn 0 trong không gian hiện hành.}
\ucflow{Luồng thực thi chính}
\ucstep{1}{Người dùng}{Trên màn hình Tủ lạnh, chọn một thực phẩm và nhấn “Đã dùng”.}
\ucstep{2}{Hệ thống}{Hiển thị tổng lượng còn lại, các lô theo thứ tự ngày hết hạn tăng dần và ô nhập lượng sử dụng với các nút nhanh “Một nửa”, “Hết”.}
\ucstep{3}{Người dùng}{Nhập lượng sử dụng hoặc chọn nút nhanh, nhấn “Xác nhận”.}
\ucstep{4}{Hệ thống}{Kiểm tra lượng sử dụng dương và không vượt tổng lượng còn lại của các lô chưa hết hạn (\br{BR14}).}
\ucstep{5}{Hệ thống}{Trừ dần từ lô có ngày hết hạn sớm nhất (\br{BR15}); lô có lượng còn lại bằng 0 chuyển sang Đã dùng hết.}
\ucstep{6}{Hệ thống}{Ghi nhật ký tiêu thụ cho từng lô bị trừ, cập nhật hiển thị và thông báo \msg{MSG10}.}
\ucflow{Luồng thực thi mở rộng}
\ucstep{3a}{Người dùng}{Chọn đích danh một lô (ví dụ hộp sữa đã mở): hệ thống chỉ trừ trên lô được chọn.}
\ucstep{4a}{Hệ thống}{Lượng sử dụng vượt lượng còn lại: hiển thị \msg{MSG17}; quay lại bước 3.}
\ucstep{4b}{Hệ thống}{Chỉ còn lô Đã hết hạn: hỏi người dùng có chắc chắn sử dụng lô quá hạn hay muốn chuyển sang \uc{UC29} Loại bỏ thực phẩm.}
\ucfield{Điều kiện sau}{Tồn kho của thực phẩm giảm đúng bằng lượng sử dụng; nhật ký tiêu thụ được ghi nhận để phục vụ \uc{UC45}.}
\ucfield{Quy tắc nghiệp vụ}{\br{BR14}, \br{BR15}, \br{BR16}}
\end{ucspec}

\uc{UC31} thực hiện yêu cầu tự động nhắc trước hạn sử dụng. Hệ thống tuân theo ba nguyên tắc: áp dụng thiết lập của từng thành viên, gộp nhiều thực phẩm vào một thông báo và không nhắc cùng một lô quá một lần trong ngày. Luồng xử lý và tương tác giữa các thành phần được thể hiện lần lượt ở Hình~\ref{fig:act-nhac-nho} và Hình~\ref{fig:seq-nhac-nho}.

\begin{ucspec}{UC31}{Gửi nhắc nhở thực phẩm sắp hết hạn}
\ucfield{Tác nhân}{Bộ lập lịch (tác nhân thời gian, khởi phát use case); Dịch vụ thông báo đẩy (phụ)}
\ucfield{Mô tả}{Hằng ngày, hệ thống tự động cập nhật trạng thái hạn dùng của mọi lô thực phẩm và gửi thông báo tổng hợp tới từng thành viên về các thực phẩm sắp hết hạn trong vòng N ngày.}
\ucfield{Điều kiện trước}{Đến khung giờ chạy tác vụ nhắc nhở được cấu hình tại \uc{UC53} (mặc định 08:00).}
\ucflow{Luồng thực thi chính}
\ucstep{1}{Bộ lập lịch}{Kích hoạt tác vụ kiểm tra hạn dùng.}
\ucstep{2}{Hệ thống}{Tính lại số ngày còn lại và trạng thái hạn dùng cho mọi lô có lượng còn lại lớn hơn 0 (\br{BR12}).}
\ucstep{3}{Hệ thống}{Xác định các không gian có lô còn không quá 7 ngày, là giá trị lớn nhất của ngưỡng N, hoặc vừa chuyển sang Đã hết hạn trong ngày.}
\ucstep{4}{Hệ thống}{Với mỗi thành viên của các không gian đó, đọc thiết lập nhắc nhở cá nhân (\uc{UC08}) gồm ngưỡng N và trạng thái bật, tắt.}
\ucstep{5}{Hệ thống}{Lọc các lô có số ngày còn lại từ 0 đến N, hoặc vừa quá hạn trong ngày, và chưa được nhắc cho thành viên này trong ngày (\br{BR13}).}
\ucstep{6}{Hệ thống}{Tạo một thông báo tổng hợp, ví dụ “3 thực phẩm sắp hết hạn: Thịt ba chỉ (còn 1 ngày), Sữa tươi (còn 2 ngày)…”, liệt kê tối đa 5 lô gần hết hạn nhất, lưu vào hộp thư trong ứng dụng.}
\ucstep{7}{Hệ thống}{Gửi yêu cầu tới Dịch vụ thông báo đẩy cho mọi thiết bị đã đăng ký của thành viên.}
\ucstep{8}{Dịch vụ thông báo đẩy}{Chuyển thông báo tới thiết bị; khi người dùng chạm vào, ứng dụng mở màn hình \uc{UC30} Xem thực phẩm sắp hết hạn kèm món gợi ý.}
\ucflow{Luồng thực thi mở rộng}
\ucstep{1a}{Hệ thống}{Lần chạy trước bị gián đoạn: tác vụ được thiết kế lũy đẳng, chạy lại không gây gửi trùng nhờ bản ghi đã nhắc theo ngày.}
\ucstep{4a}{Hệ thống}{Thành viên tắt nhắc nhở hạn dùng: bỏ qua thành viên này.}
\ucstep{5a}{Hệ thống}{Không còn lô nào cần nhắc cho thành viên: không tạo thông báo.}
\ucstep{7a}{Hệ thống}{Mã thiết bị không còn hợp lệ: xóa mã thiết bị khỏi tài khoản.}
\ucstep{7b}{Hệ thống}{Dịch vụ thông báo đẩy lỗi tạm thời: thử lại tối đa 3 lần với khoảng chờ tăng dần, ghi nhật ký nếu vẫn thất bại; thông báo vẫn có trong hộp thư (\nfr{NFR12}).}
\ucfield{Điều kiện sau}{Mọi lô có lượng còn lại được cập nhật trạng thái hạn dùng; mỗi thành viên bật nhắc nhở nhận tối đa một thông báo tổng hợp mỗi ngày cho các lô thuộc ngưỡng của mình.}
\ucfield{Quy tắc nghiệp vụ}{\br{BR12}, \br{BR13}, \br{BR23}}
\end{ucspec}

\diagram[0.64\textwidth]{c2-11-act-nhac-nho}{Biểu đồ hoạt động của use case Gửi nhắc nhở thực phẩm sắp hết hạn}{fig:act-nhac-nho}

\diagram[\textwidth]{c2-12-seq-nhac-nho}{Biểu đồ tuần tự của luồng nhắc nhở hạn dùng}{fig:seq-nhac-nho}

Theo biểu đồ tuần tự ở Hình~\ref{fig:seq-nhac-nho}, thông báo đẩy chỉ hiển thị nội dung tóm tắt. Khi người dùng mở thông báo, ứng dụng mới truy vấn danh sách chi tiết và món gợi ý từ máy chủ. Cách này hạn chế thông tin hiển thị trên màn hình khóa và giúp ứng dụng tải dữ liệu mới nhất. Các use case \uc{UC24}, \uc{UC26}, \uc{UC27}, \uc{UC29} và \uc{UC30} được đặc tả tại Phụ lục~\ref{pl:dac-ta-bo-sung}.

% ---------------------------------------------------------------------
%  2.2.2.5. Phân rã gói "Quản lý công thức nấu ăn"
% ---------------------------------------------------------------------
\subsubsection{Phân rã use case “Quản lý công thức nấu ăn”}\label{sssec:pr-cong-thuc}

Gói Quản lý công thức nấu ăn gồm năm use case, được phân rã tại Hình~\ref{fig:uc-cong-thuc}. Ngoại trừ \uc{UC32} Thêm công thức, các thao tác khác bắt đầu từ màn hình chi tiết; vì vậy \uc{UC33}, \uc{UC35} và \uc{UC36} mở rộng \uc{UC34}. Use case \uc{UC36} liên kết công thức với danh sách mua sắm và chỉ khả dụng khi còn thiếu nguyên liệu.

\diagram[0.78\textwidth]{c2-13-uc-cong-thuc}{Biểu đồ use case phân rã gói Quản lý công thức nấu ăn}{fig:uc-cong-thuc}

\begin{ucspec}{UC32}{Thêm công thức nấu ăn}
\ucfield{Tác nhân}{Người dùng}
\ucfield{Mô tả}{Cho phép người dùng lưu một công thức nấu ăn vào bộ sưu tập cá nhân, với nguyên liệu được định lượng và liên kết tới danh mục thực phẩm.}
\ucfield{Điều kiện trước}{Người dùng đã đăng nhập.}
\ucflow{Luồng thực thi chính}
\ucstep{1}{Người dùng}{Tại màn hình Công thức, nhấn “Thêm công thức”.}
\ucstep{2}{Hệ thống}{Hiển thị biểu mẫu nhiều bước: thông tin chung, nguyên liệu, các bước thực hiện (Bảng~\ref{tab:in-cong-thuc}).}
\ucstep{3}{Người dùng}{Nhập tên món, chọn ảnh, số khẩu phần chuẩn, thời gian chuẩn bị và nấu, bữa phù hợp.}
\ucstep{4}{Người dùng}{Thêm từng nguyên liệu bằng cách chọn thực phẩm từ gợi ý, nhập số lượng, đơn vị và đánh dấu bắt buộc hay tùy chọn.}
\ucstep{5}{Người dùng}{Nhập các bước thực hiện theo thứ tự, có thể kéo thả để sắp xếp lại; nhấn “Lưu”.}
\ucstep{6}{Hệ thống}{Kiểm tra các trường bắt buộc, độ dài, miền giá trị và yêu cầu có ít nhất một nguyên liệu (\br{BR17}).}
\ucstep{7}{Hệ thống}{Lưu công thức vào bộ sưu tập của người dùng và hiển thị màn hình chi tiết công thức (\uc{UC34}).}
\ucflow{Luồng thực thi mở rộng}
\ucstep{4a}{Người dùng}{Nguyên liệu chưa có trong danh mục: tạo thực phẩm riêng như ở luồng 3a của \uc{UC25}.}
\ucstep{4b}{Hệ thống}{Cùng một thực phẩm được thêm hai lần: gộp số lượng theo \br{BR16} và thông báo cho người dùng.}
\ucstep{5a}{Người dùng}{Thoát giữa chừng: hệ thống lưu bản nháp trên thiết bị và đề nghị khôi phục ở lần mở biểu mẫu kế tiếp.}
\ucstep{6a}{Hệ thống}{Dữ liệu không hợp lệ: chuyển về bước có lỗi, hiển thị \msg{MSG01} hoặc \msg{MSG02} tại trường tương ứng.}
\ucfield{Điều kiện sau}{Công thức mới thuộc bộ sưu tập của người dùng, sẵn sàng được đối chiếu với tủ lạnh và tham gia gợi ý món (\uc{UC40}).}
\ucfield{Quy tắc nghiệp vụ}{\br{BR16}, \br{BR17}}
\end{ucspec}

\begin{fieldtable}{Dữ liệu đầu vào của một công thức nấu ăn}{tab:in-cong-thuc}
\frow{Tên món}{Tên hiển thị của công thức}{Có}{Dài 2–100 ký tự}{Canh chua cá lóc}
\frow{Ảnh minh họa}{Ảnh món ăn}{Không}{Định dạng JPEG hoặc PNG, tối đa 5 MB, tự động nén}{canh-chua.jpg}
\frow{Số khẩu phần chuẩn}{Số người ăn ứng với định lượng}{Có}{Số nguyên từ 1 đến 20}{4}
\frow{Thời gian}{Thời gian chuẩn bị và thời gian nấu}{Không}{Số nguyên phút, từ 0 đến 600}{15 và 25}
\frow{Bữa phù hợp}{Bữa thường dùng món này}{Không}{Chọn một hoặc nhiều trong sáng, trưa, tối}{Trưa, tối}
\frow{Nguyên liệu}{Danh sách thực phẩm, số lượng, đơn vị, bắt buộc hay tùy chọn}{Có}{Ít nhất 1 nguyên liệu; số lượng dương}{Cá lóc 500 g (bắt buộc)}
\frow{Các bước thực hiện}{Hướng dẫn nấu theo thứ tự}{Không}{Mỗi bước tối đa 500 ký tự}{Bước 1: Sơ chế cá…}
\frow{Ghi chú}{Mẹo nấu, lưu ý dinh dưỡng}{Không}{Tối đa 1000 ký tự}{Có thể thay me bằng sấu}
\end{fieldtable}

\begin{ucspec}{UC34}{Xem chi tiết công thức và đối chiếu nguyên liệu}
\ucfield{Tác nhân}{Người dùng}
\ucfield{Mô tả}{Cho phép người dùng xem đầy đủ một công thức, thay đổi số khẩu phần và biết ngay nguyên liệu nào đã có, nguyên liệu nào còn thiếu trong tủ lạnh.}
\ucfield{Điều kiện trước}{Người dùng chọn một công thức từ bộ sưu tập, từ công thức mẫu, từ kết quả tìm kiếm hoặc từ danh sách gợi ý.}
\ucflow{Luồng thực thi chính}
\ucstep{1}{Người dùng}{Chọn một công thức.}
\ucstep{2}{Hệ thống}{Hiển thị ảnh, tên, thời gian, số khẩu phần, các bước thực hiện và danh sách nguyên liệu.}
\ucstep{3}{Hệ thống}{Với mỗi nguyên liệu, so sánh lượng cần với tổng lượng chưa hết hạn trong tủ của không gian hiện hành sau khi quy đổi đơn vị (\br{BR16}); gắn nhãn “Đủ”, “Thiếu x” hoặc “Không có”, và đánh dấu nguyên liệu đang có lô sắp hết hạn.}
\ucstep{4}{Người dùng}{Thay đổi số khẩu phần bằng nút tăng, giảm.}
\ucstep{5}{Hệ thống}{Nhân định lượng theo tỷ lệ mới và đối chiếu lại như bước 3.}
\ucflow{Luồng thực thi mở rộng}
\ucstep{5a}{Người dùng}{Có nguyên liệu thiếu và chọn “Thêm phần thiếu vào danh sách mua sắm”: thực hiện \uc{UC36}.}
\ucstep{5b}{Người dùng}{Chọn biểu tượng trái tim: thực hiện \uc{UC35}.}
\ucstep{5c}{Người dùng}{Công thức thuộc bộ sưu tập của mình và chọn “Sửa” hoặc “Xóa”: thực hiện \uc{UC33}.}
\ucstep{5d}{Người dùng}{Công thức là công thức mẫu và chọn “Lưu về bộ sưu tập”: hệ thống tạo bản sao riêng có thể chỉnh sửa (\br{BR17}).}
\ucstep{5e}{Người dùng}{Chọn “Thêm vào lịch”: thực hiện \uc{UC37} với món đã được chọn sẵn.}
\ucfield{Điều kiện sau}{Không thay đổi dữ liệu, trừ khi người dùng thực hiện một trong các use case mở rộng.}
\ucfield{Quy tắc nghiệp vụ}{\br{BR15}, \br{BR16}, \br{BR17}}
\end{ucspec}

Các use case \uc{UC33}, \uc{UC35} và \uc{UC36} được đặc tả tại Phụ lục~\ref{pl:dac-ta-bo-sung}.

% ---------------------------------------------------------------------
%  2.2.2.6. Phân rã gói "Quản lý kế hoạch bữa ăn"
% ---------------------------------------------------------------------
\subsubsection{Phân rã use case “Quản lý kế hoạch bữa ăn”}\label{sssec:pr-ke-hoach}

Gói Quản lý kế hoạch bữa ăn hiện thực chức năng tương ứng trong đề bài và được phân rã tại Hình~\ref{fig:uc-ke-hoach}. Hai điểm vào chính là \uc{UC37} Lên lịch bữa ăn và \uc{UC38} Xem lịch bữa ăn. Use case \uc{UC40} Gợi ý món ăn hoạt động độc lập, đồng thời mở rộng \uc{UC37} như một cách chọn món. Từ màn hình lịch, người dùng cũng có thể cập nhật lịch (\uc{UC39}), xác nhận đã nấu (\uc{UC41}) hoặc sinh danh sách mua sắm (\uc{UC22}). \uc{UC22} thuộc gói mua sắm nhưng được thể hiện tại đây để làm rõ mối liên hệ giữa hai gói.

\diagram[0.82\textwidth]{c2-14-uc-ke-hoach}{Biểu đồ use case phân rã gói Quản lý kế hoạch bữa ăn}{fig:uc-ke-hoach}

\begin{ucspec}{UC37}{Lên lịch bữa ăn}
\ucfield{Tác nhân}{Người dùng}
\ucfield{Mô tả}{Cho phép người dùng thêm một hoặc nhiều món vào bữa sáng, trưa hoặc tối của một ngày trong kế hoạch.}
\ucfield{Điều kiện trước}{Người dùng đã đăng nhập; có ít nhất một công thức trong bộ sưu tập hoặc trong công thức mẫu.}
\ucflow{Luồng thực thi chính}
\ucstep{1}{Người dùng}{Tại màn hình Kế hoạch, chọn một ngày và nhấn “+” tại bữa cần lên lịch.}
\ucstep{2}{Hệ thống}{Hiển thị ba thẻ chọn món: “Gợi ý cho bạn”, “Yêu thích”, “Tất cả công thức”, mặc định mở thẻ gợi ý.}
\ucstep{3}{Người dùng}{Chọn một hoặc nhiều món, nhập số khẩu phần cho mỗi món (mặc định bằng số thành viên của nhóm) và nhấn “Thêm”.}
\ucstep{4}{Hệ thống}{Kiểm tra ngày nằm trong khoảng từ hôm nay đến 28 ngày tiếp theo và số khẩu phần từ 1 đến 20 (\br{BR18}).}
\ucstep{5}{Hệ thống}{Thêm các món vào bữa ở trạng thái Dự kiến, đối chiếu nguyên liệu và hiển thị biểu tượng đủ hoặc thiếu nguyên liệu trên từng món.}
\ucflow{Luồng thực thi mở rộng}
\ucstep{2a}{Người dùng}{Ở thẻ “Gợi ý cho bạn”: thực hiện \uc{UC40} với bữa và số khẩu phần đang chọn.}
\ucstep{3a}{Người dùng}{Nhập tên món tự do không có công thức (ví dụ “Ăn ngoài”): món được lưu nhưng không tham gia đối chiếu nguyên liệu.}
\ucstep{4a}{Hệ thống}{Ngày đã qua hoặc vượt 28 ngày: hiển thị \msg{MSG28}; quay lại bước 1.}
\ucstep{5a}{Hệ thống}{Món đã có trong cùng bữa: hỏi người dùng muốn tăng số khẩu phần của món hiện có hay thêm một dòng mới.}
\ucfield{Điều kiện sau}{Bữa ăn của ngày được chọn chứa các món mới ở trạng thái Dự kiến, kèm thông tin đủ hoặc thiếu nguyên liệu.}
\ucfield{Quy tắc nghiệp vụ}{\br{BR18}}
\end{ucspec}

Use case Gợi ý món ăn áp dụng cơ chế chấm điểm tại mục~\ref{ssec:nv05}, theo các công thức~\eqref{eq:diem-han-dung} và~\eqref{eq:diem-goi-y}. Đặc tả và luồng tính toán được trình bày lần lượt tại Bảng~\ref{tab:UC40} và Hình~\ref{fig:act-goi-y}.

\begin{ucspec}{UC40}{Gợi ý món ăn từ thực phẩm sẵn có}
\ucfield{Tác nhân}{Người dùng}
\ucfield{Mô tả}{Hệ thống xếp hạng các công thức theo mức độ đáp ứng của tủ lạnh hiện tại, ưu tiên những món giúp sử dụng thực phẩm sắp hết hạn, và đề xuất tối đa 10 món.}
\ucfield{Điều kiện trước}{Người dùng đã đăng nhập; người dùng mở chức năng gợi ý từ màn hình chính, từ màn hình thực phẩm sắp hết hạn hoặc từ \uc{UC37}.}
\ucflow{Luồng thực thi chính}
\ucstep{1}{Người dùng}{Mở “Gợi ý món ăn”, tùy chọn bữa và số khẩu phần.}
\ucstep{2}{Hệ thống}{Lấy tồn kho khả dụng gồm các lô Còn hạn và Sắp hết hạn của không gian hiện hành.}
\ucstep{3}{Hệ thống}{Lấy tập công thức ứng viên gồm bộ sưu tập của người dùng và công thức mẫu, lọc theo bữa phù hợp nếu người dùng đã chọn bữa.}
\ucstep{4}{Hệ thống}{Với mỗi công thức, tính tỷ lệ đáp ứng $R$, điểm ưu tiên hạn dùng $E$ và hệ số yêu thích $F$; loại các công thức có $R < \theta$ (\br{BR19}).}
\ucstep{5}{Hệ thống}{Tính điểm $S = 0{,}6R + 0{,}3E + 0{,}1F$, sắp xếp giảm dần và lấy tối đa 10 món.}
\ucstep{6}{Hệ thống}{Hiển thị mỗi món kèm nhãn “Đủ nguyên liệu” hoặc “Thiếu k nguyên liệu”, danh sách thực phẩm sắp hết hạn mà món sử dụng và thời gian nấu.}
\ucstep{7}{Người dùng}{Chọn một món và chọn “Thêm vào lịch” hoặc “Xem công thức”.}
\ucflow{Luồng thực thi mở rộng}
\ucstep{2a}{Hệ thống}{Tủ lạnh trống: hiển thị các công thức mẫu phổ biến và lời nhắc thêm thực phẩm (\uc{UC25}).}
\ucstep{4a}{Hệ thống}{Không công thức nào đạt ngưỡng: hiển thị \msg{MSG31} cùng ba món có $R$ cao nhất và gợi ý bổ sung nguyên liệu còn thiếu.}
\ucstep{7a}{Người dùng}{Chọn “Thêm vào lịch”: thực hiện bước 3 của \uc{UC37} với món đã chọn.}
\ucstep{7b}{Người dùng}{Chọn “Xem công thức”: thực hiện \uc{UC34}.}
\ucfield{Điều kiện sau}{Không thay đổi dữ liệu; kết quả gợi ý được lưu đệm trong 10 phút hoặc đến khi tồn kho thay đổi.}
\ucfield{Quy tắc nghiệp vụ}{\br{BR12}, \br{BR16}, \br{BR19}}
\end{ucspec}

\diagram[0.64\textwidth]{c2-15-act-goi-y-mon}{Biểu đồ hoạt động của use case Gợi ý món ăn từ thực phẩm sẵn có}{fig:act-goi-y}

Xét ví dụ một gia đình có cá lóc còn 1 ngày, dứa còn 2 ngày, cà chua còn 5 ngày và trứng gà còn 10 ngày trước khi hết hạn; ngưỡng nhắc là $N = 3$. Tủ lạnh cũng có các gia vị cơ bản. Công thức canh chua cá lóc cần bốn nguyên liệu bắt buộc: cá lóc, cà chua, dứa và me. Do thiếu me, tỷ lệ đáp ứng là $R = 3/4 = 0{,}75$. Trọng số hạn dùng của cá lóc là $(3 + 1 - 1)/4 = 0{,}75$, của dứa là $(3 + 1 - 2)/4 = 0{,}5$, còn cà chua nằm ngoài ngưỡng nên có trọng số 0. Vì vậy, $E = \min(1;\ 0{,}75 + 0{,}5) = 1$. Nếu chưa được đánh dấu yêu thích, món này có $F = 0$ và $S = 0{,}6 \times 0{,}75 + 0{,}3 \times 1 = 0{,}75$. Món trứng xào cà chua có đủ nguyên liệu nhưng không dùng thực phẩm sắp hết hạn, nên $R = 1$, $E = 0$ và $S = 0{,}6$. Do đó, canh chua cá lóc được xếp hạng cao hơn và hệ thống xác định me là nguyên liệu cần mua.

\begin{ucspec}{UC41}{Xác nhận đã nấu món ăn}
\ucfield{Tác nhân}{Người dùng}
\ucfield{Mô tả}{Cho phép người dùng xác nhận một món trong lịch đã được nấu; hệ thống trừ nguyên liệu tương ứng khỏi tủ lạnh theo nguyên tắc FEFO và ghi nhận tiêu thụ.}
\ucfield{Điều kiện trước}{Món ở trạng thái Dự kiến, thuộc ngày hôm nay hoặc ngày đã qua, và gắn với một công thức.}
\ucflow{Luồng thực thi chính}
\ucstep{1}{Người dùng}{Trên lịch bữa ăn, chọn món và nhấn “Đã nấu”.}
\ucstep{2}{Hệ thống}{Tính lượng nguyên liệu cần dùng theo số khẩu phần thực tế (mặc định bằng số khẩu phần dự kiến).}
\ucstep{3}{Hệ thống}{Với mỗi nguyên liệu, lấy các lô chưa hết hạn (Còn hạn và Sắp hết hạn) theo thứ tự ngày hết hạn tăng dần và phân bổ lượng trừ từ lô sớm nhất (\br{BR15}).}
\ucstep{4}{Hệ thống}{Hiển thị bảng xem trước: nguyên liệu, lượng sẽ trừ ở từng lô, lượng thiếu nếu tồn kho không đủ.}
\ucstep{5}{Người dùng}{Điều chỉnh lượng thực dùng (nếu dùng ít hoặc nhiều hơn công thức) và nhấn “Xác nhận”.}
\ucstep{6}{Hệ thống}{Cập nhật lượng còn lại của các lô trong một giao dịch duy nhất; lô về 0 chuyển sang Đã dùng hết (\br{BR14}).}
\ucstep{7}{Hệ thống}{Ghi nhật ký tiêu thụ gắn với món ăn, chuyển món sang trạng thái Đã nấu và thông báo \msg{MSG10}.}
\ucflow{Luồng thực thi mở rộng}
\ucstep{4a}{Hệ thống}{Có nguyên liệu không đủ hoặc không có trong tủ: đánh dấu dòng thiếu; người dùng có thể bỏ qua (coi như đã dùng nguồn ngoài tủ) hoặc bổ sung tồn kho trước.}
\ucstep{4b}{Hệ thống}{Chỉ còn lô Đã hết hạn cho một nguyên liệu: không tự động trừ, yêu cầu người dùng chọn rõ theo \br{BR15}.}
\ucstep{5a}{Người dùng}{Chọn “Bỏ qua trừ kho”: món chuyển sang Đã nấu mà không thay đổi tồn kho.}
\ucstep{6a}{Hệ thống}{Có thay đổi tồn kho đồng thời từ thành viên khác làm lượng còn lại không đủ: hủy giao dịch, tải lại số liệu và quay lại bước 4.}
\ucfield{Điều kiện sau}{Món ở trạng thái Đã nấu và không thể sửa (\br{BR18}); tồn kho giảm đúng lượng đã xác nhận; nhật ký tiêu thụ được ghi nhận.}
\ucfield{Quy tắc nghiệp vụ}{\br{BR14}, \br{BR15}, \br{BR16}, \br{BR18}}
\end{ucspec}

\diagram[0.6\textwidth]{c2-16-act-xac-nhan-nau}{Biểu đồ hoạt động của use case Xác nhận đã nấu món ăn}{fig:act-xac-nhan-nau}

Bước 6 của \uc{UC41} chạy trong một giao dịch duy nhất: hệ thống cập nhật toàn bộ lô liên quan hoặc không cập nhật lô nào. Cách xử lý này bảo đảm tính nhất quán khi nhiều thành viên cùng thao tác trên một tủ lạnh. Các use case \uc{UC38} và \uc{UC39} được đặc tả tại Phụ lục~\ref{pl:dac-ta-bo-sung}.

% ---------------------------------------------------------------------
%  2.2.2.7. Phân rã gói "Tìm kiếm và phân loại thực phẩm"
% ---------------------------------------------------------------------
\subsubsection{Phân rã use case “Tìm kiếm và phân loại thực phẩm”}\label{sssec:pr-tim-kiem}

Gói Tìm kiếm và phân loại thực phẩm gồm hai use case: \uc{UC42} và \uc{UC43}. Hình~\ref{fig:uc-tim-kiem} còn thể hiện ba use case thuộc các gói khác để cho thấy các bước có thể thực hiện từ kết quả tìm kiếm. Người dùng có thể mở chi tiết công thức (\uc{UC34}) hoặc thêm thực phẩm vào tủ (\uc{UC25}). Use case \uc{UC40} cũng xuất hiện ở đây theo cách phân nhóm của đề bài, dù cơ chế gợi ý món được dùng chung với gói kế hoạch bữa ăn.

\diagram[0.82\textwidth]{c2-17-uc-tim-kiem}{Biểu đồ use case phân rã gói Tìm kiếm và phân loại thực phẩm}{fig:uc-tim-kiem}

\begin{ucspec}{UC42}{Tìm kiếm thực phẩm, món ăn}
\ucfield{Tác nhân}{Người dùng}
\ucfield{Mô tả}{Cho phép người dùng tìm thực phẩm, món ăn trong danh mục hệ thống, trong tủ lạnh của không gian hiện hành hoặc trong bộ sưu tập công thức bằng từ khóa có dấu hoặc không dấu.}
\ucfield{Điều kiện trước}{Người dùng đã đăng nhập.}
\ucflow{Luồng thực thi chính}
\ucstep{1}{Người dùng}{Chạm vào thanh tìm kiếm ở đầu màn hình chính.}
\ucstep{2}{Hệ thống}{Hiển thị các từ khóa tìm gần đây và các loại thực phẩm phổ biến.}
\ucstep{3}{Người dùng}{Gõ từ khóa, ví dụ “ca rot”.}
\ucstep{4}{Hệ thống}{Sau 300 ms kể từ lần gõ cuối, chuẩn hóa từ khóa về chữ thường không dấu và tìm đồng thời trong ba nguồn: tủ lạnh, công thức, danh mục thực phẩm.}
\ucstep{5}{Hệ thống}{Hiển thị kết quả theo ba nhóm “Trong tủ lạnh”, “Món ăn”, “Thực phẩm”, xếp hạng theo mức khớp: trùng khớp hoàn toàn, khớp tiền tố, khớp một phần; mỗi nhóm hiển thị tối đa 5 kết quả kèm nút “Xem thêm”.}
\ucstep{6}{Người dùng}{Chọn một kết quả.}
\ucstep{7}{Hệ thống}{Mở màn hình tương ứng: chi tiết lô trong tủ (\uc{UC24}), chi tiết công thức (\uc{UC34}) hoặc biểu mẫu thêm thực phẩm đã điền sẵn (\uc{UC25}).}
\ucflow{Luồng thực thi mở rộng}
\ucstep{4a}{Hệ thống}{Mất kết nối mạng: chỉ tìm trong dữ liệu tủ lạnh và công thức đã lưu đệm trên thiết bị, ghi chú kết quả có thể chưa đầy đủ.}
\ucstep{5a}{Hệ thống}{Không có kết quả: hiển thị \msg{MSG18} và gợi ý “Thêm ‘<từ khóa>’ vào danh sách mua sắm”.}
\ucstep{5b}{Người dùng}{Chạm biểu tượng bộ lọc: thực hiện \uc{UC43} trên tập kết quả hiện tại.}
\ucfield{Điều kiện sau}{Không thay đổi dữ liệu; từ khóa được lưu vào lịch sử tìm kiếm trên thiết bị.}
\ucfield{Quy tắc nghiệp vụ}{\br{BR23}}
\end{ucspec}

Use case \uc{UC43} Lọc thực phẩm theo danh mục được đặc tả tại Phụ lục~\ref{pl:dac-ta-bo-sung}.

% ---------------------------------------------------------------------
%  2.2.2.8. Phân rã gói "Thống kê báo cáo"
% ---------------------------------------------------------------------
\subsubsection{Phân rã use case “Thống kê báo cáo”}\label{sssec:pr-thong-ke}

Gói Thống kê báo cáo gồm ba loại báo cáo và chức năng xuất dữ liệu, được phân rã tại Hình~\ref{fig:uc-thong-ke}. Ba use case \uc{UC44}, \uc{UC45} và \uc{UC46} dùng chung bộ chọn thời gian, các thẻ chỉ số và biểu đồ chi tiết, nhưng khai thác các nguồn dữ liệu khác nhau (Bảng~\ref{tab:chi-so-bao-cao}). Use case \uc{UC47} cho phép xuất từng báo cáo. Số liệu được tính theo không gian hiện hành; trong không gian gia đình, báo cáo phản ánh hoạt động chung của các thành viên.

\diagram[0.62\textwidth]{c2-18-uc-thong-ke}{Biểu đồ use case phân rã gói Thống kê báo cáo}{fig:uc-thong-ke}

\begin{ucspec}{UC44}{Xem báo cáo mua sắm và chi tiêu}
\ucfield{Tác nhân}{Người dùng}
\ucfield{Mô tả}{Cho phép người dùng xem tổng chi tiêu, cơ cấu chi tiêu theo loại thực phẩm, xu hướng theo thời gian và các mặt hàng mua nhiều nhất trong một khoảng thời gian.}
\ucfield{Điều kiện trước}{Người dùng đã đăng nhập.}
\ucflow{Luồng thực thi chính}
\ucstep{1}{Người dùng}{Mở màn hình Thống kê, chọn thẻ “Mua sắm”.}
\ucstep{2}{Hệ thống}{Hiển thị bộ chọn khoảng thời gian với các lựa chọn nhanh: tuần này, tháng này, 3 tháng gần nhất, tùy chọn; mặc định là tháng này.}
\ucstep{3}{Người dùng}{Chọn khoảng thời gian và cách nhóm theo ngày, tuần hoặc tháng.}
\ucstep{4}{Hệ thống}{Kiểm tra khoảng thời gian hợp lệ, không quá 12 tháng (\br{BR21}).}
\ucstep{5}{Hệ thống}{Tổng hợp từ các mặt hàng đã mua, kể cả mặt hàng đã chuyển sang Đã nhập tủ lạnh, trong khoảng: tổng chi tiêu, số lần đi chợ, chi tiêu trung bình mỗi lần, so sánh với kỳ liền trước.}
\ucstep{6}{Hệ thống}{Hiển thị các thẻ chỉ số, biểu đồ cột chi tiêu theo thời gian, biểu đồ tròn cơ cấu theo loại và bảng mười mặt hàng mua nhiều nhất.}
\ucstep{7}{Người dùng}{Chạm vào một cột hoặc một phần của biểu đồ để xem chi tiết các mặt hàng tương ứng.}
\ucflow{Luồng thực thi mở rộng}
\ucstep{4a}{Hệ thống}{Khoảng thời gian không hợp lệ: hiển thị \msg{MSG25}; quay lại bước 3.}
\ucstep{5a}{Hệ thống}{Không có dữ liệu trong khoảng: hiển thị \msg{MSG26} cùng hướng dẫn ghi đơn giá khi đánh dấu đã mua.}
\ucstep{5b}{Hệ thống}{Có mặt hàng không ghi đơn giá: tính vào số lượng nhưng không tính vào chi tiêu, hiển thị ghi chú “k mặt hàng chưa có giá”.}
\ucstep{7a}{Người dùng}{Chọn “Xuất báo cáo”: thực hiện \uc{UC47}.}
\ucfield{Điều kiện sau}{Không thay đổi dữ liệu.}
\ucfield{Quy tắc nghiệp vụ}{\br{BR21}, \br{BR23}}
\end{ucspec}

Báo cáo lãng phí tổng hợp giá trị thực phẩm bị loại bỏ và các mặt hàng thường bị bỏ đi. Người dùng có thể dựa vào số liệu này để điều chỉnh lượng mua.

\begin{ucspec}{UC46}{Xem báo cáo lãng phí thực phẩm}
\ucfield{Tác nhân}{Người dùng}
\ucfield{Mô tả}{Cho phép người dùng xem số lượng, giá trị ước tính và tỷ lệ thực phẩm bị loại bỏ do hỏng hoặc hết hạn, cùng danh sách thực phẩm hay bị lãng phí.}
\ucfield{Điều kiện trước}{Người dùng đã đăng nhập.}
\ucflow{Luồng thực thi chính}
\ucstep{1}{Người dùng}{Mở màn hình Thống kê, chọn thẻ “Lãng phí”.}
\ucstep{2}{Hệ thống}{Hiển thị bộ chọn khoảng thời gian như bước 2 của \uc{UC44}.}
\ucstep{3}{Người dùng}{Chọn khoảng thời gian.}
\ucstep{4}{Hệ thống}{Kiểm tra khoảng thời gian không quá 12 tháng (\br{BR21}).}
\ucstep{5}{Hệ thống}{Tổng hợp nhật ký loại bỏ: số lô bị loại bỏ, giá trị lãng phí theo đơn giá của mặt hàng nguồn, tỷ lệ lãng phí theo Bảng~\ref{tab:chi-so-bao-cao}, phân bổ theo lý do (hết hạn, hỏng, khác).}
\ucstep{6}{Hệ thống}{Hiển thị các thẻ chỉ số kèm mức thay đổi so với kỳ trước, biểu đồ xu hướng theo tuần và danh sách năm thực phẩm bị loại bỏ nhiều nhất cùng gợi ý giảm lượng mua.}
\ucflow{Luồng thực thi mở rộng}
\ucstep{4a}{Hệ thống}{Khoảng thời gian không hợp lệ: hiển thị \msg{MSG25}; quay lại bước 3.}
\ucstep{5a}{Hệ thống}{Không có thực phẩm nào bị loại bỏ: hiển thị thông điệp khích lệ thay cho biểu đồ.}
\ucstep{6a}{Người dùng}{Chọn “Xuất báo cáo”: thực hiện \uc{UC47}.}
\ucfield{Điều kiện sau}{Không thay đổi dữ liệu.}
\ucfield{Quy tắc nghiệp vụ}{\br{BR21}, \br{BR23}}
\end{ucspec}

Các use case \uc{UC45} và \uc{UC47} được đặc tả tại Phụ lục~\ref{pl:dac-ta-bo-sung}.

% ---------------------------------------------------------------------
%  2.2.2.9. Phân rã gói "Quản trị hệ thống"
% ---------------------------------------------------------------------
\subsubsection{Phân rã use case “Quản trị hệ thống”}\label{sssec:pr-quan-tri}

Gói Quản trị hệ thống gồm sáu use case dành cho Quản trị viên, được phân rã tại Hình~\ref{fig:uc-quan-tri}. Use case \uc{UC48} quản lý tài khoản, \uc{UC53} cấu hình tham số và bốn use case còn lại quản lý dữ liệu danh mục dùng chung. \uc{UC49} Quản lý danh mục thực phẩm bao gồm \uc{UC50} và \uc{UC51}, vì mỗi thực phẩm phải thuộc một loại và có đơn vị mặc định. Các chức năng này nằm trong khu vực riêng của ứng dụng di động, chỉ tài khoản quản trị mới truy cập được.

\diagram[0.82\textwidth]{c2-19-uc-quan-tri}{Biểu đồ use case phân rã gói Quản trị hệ thống}{fig:uc-quan-tri}

\begin{ucspec}{UC49}{Quản lý danh mục thực phẩm}
\ucfield{Tác nhân}{Quản trị viên}
\ucfield{Mô tả}{Cho phép quản trị viên thêm, sửa, ngừng sử dụng các thực phẩm trong danh mục dùng chung, bao gồm hạn dùng gợi ý theo từng vị trí lưu trữ và mã vạch; đồng thời xem xét các thực phẩm riêng do người dùng đề xuất.}
\ucfield{Điều kiện trước}{Quản trị viên đã đăng nhập.}
\ucflow{Luồng thực thi chính}
\ucstep{1}{Quản trị viên}{Mở khu vực Quản trị, chọn “Danh mục thực phẩm”.}
\ucstep{2}{Hệ thống}{Hiển thị danh sách thực phẩm có phân trang, ô tìm kiếm, bộ lọc theo loại và trạng thái (Đang dùng, Ngừng sử dụng, Chờ duyệt).}
\ucstep{3}{Quản trị viên}{Nhấn “Thêm thực phẩm”.}
\ucstep{4}{Hệ thống}{Hiển thị biểu mẫu gồm tên, tên gọi khác, loại thực phẩm (\uc{UC50}), đơn vị mặc định (\uc{UC51}), mã vạch, ảnh, vị trí lưu trữ mặc định và hạn dùng gợi ý cho từng vị trí.}
\ucstep{5}{Quản trị viên}{Nhập thông tin và nhấn “Lưu”.}
\ucstep{6}{Hệ thống}{Kiểm tra trường bắt buộc, tên chưa tồn tại trong cùng loại (\br{BR22}), mã vạch chưa gán cho thực phẩm khác, hạn dùng gợi ý là số nguyên dương.}
\ucstep{7}{Hệ thống}{Lưu thực phẩm ở trạng thái Đang dùng, ghi nhật ký quản trị, hiển thị \msg{MSG10}.}
\ucflow{Luồng thực thi mở rộng}
\ucstep{3a}{Quản trị viên}{Chọn một thực phẩm và nhấn “Sửa”: hệ thống hiển thị biểu mẫu với dữ liệu hiện có; tiếp tục từ bước 5.}
\ucstep{3b}{Quản trị viên}{Chọn “Ngừng sử dụng”: hệ thống hiển thị \msg{MSG19}; khi được xác nhận, thực phẩm không còn xuất hiện trong gợi ý nhưng dữ liệu đã tham chiếu vẫn giữ nguyên (\br{BR22}).}
\ucstep{3c}{Quản trị viên}{Mở bộ lọc “Chờ duyệt”: duyệt thực phẩm riêng do người dùng đề xuất bằng cách gộp vào một thực phẩm có sẵn hoặc chuẩn hóa thành thực phẩm mới.}
\ucstep{6a}{Hệ thống}{Tên đã tồn tại trong cùng loại hoặc mã vạch đã được gán: hiển thị \msg{MSG02} kèm liên kết tới bản ghi trùng; quay lại bước 5.}
\ucstep{6b}{Hệ thống}{Thiếu trường bắt buộc: hiển thị \msg{MSG01}; quay lại bước 5.}
\ucfield{Điều kiện sau}{Danh mục thực phẩm được cập nhật và có hiệu lực ngay cho chức năng gợi ý tên, đề xuất hạn dùng và quét mã vạch của mọi người dùng.}
\ucfield{Quy tắc nghiệp vụ}{\br{BR16}, \br{BR22}}
\end{ucspec}

Use case \uc{UC53} cho phép thay đổi tham số vận hành mà không cần phát hành phiên bản mới. Bảng~\ref{tab:tham-so-he-thong} liệt kê các tham số và miền giá trị hợp lệ; giá trị mặc định tuân theo đề bài và các quy tắc nghiệp vụ.

{\small
\begin{longtable}{|L{4.6cm}|C{2.2cm}|L{3.0cm}|L{\dimexpr\textwidth-8\tabcolsep-5\arrayrulewidth-4.6cm-2.2cm-3.0cm\relax}|}
\caption{Các tham số cấu hình hệ thống}\label{tab:tham-so-he-thong}\\
\hline
\hd{Tham số} & \hd{Mặc định} & \hd{Miền hợp lệ} & \hd{Ảnh hưởng tới}\\ \hline
\endfirsthead
\multicolumn{4}{l}{\footnotesize\itshape Bảng \thetable\ (tiếp theo)}\\ \hline
\hd{Tham số} & \hd{Mặc định} & \hd{Miền hợp lệ} & \hd{Ảnh hưởng tới}\\ \hline
\endhead
Ngưỡng nhắc mặc định $N$ & 3 ngày & 1–7 ngày & Người dùng chưa tự thiết lập (\uc{UC08}, \uc{UC31})\\ \hline
Khung giờ chạy nhắc nhở & 08:00 & 05:00–21:00 & Tác vụ \uc{UC31}\\ \hline
Số thành viên tối đa mỗi nhóm & 10 & 2–30 & \uc{UC10}, \uc{UC11}\\ \hline
Thời hạn mã mời & 48 giờ & 1–168 giờ & \uc{UC10}\\ \hline
Thời hạn mã OTP & 5 phút & 2–15 phút & \uc{UC02}\\ \hline
Ngưỡng đáp ứng tối thiểu $\theta$ & 0,6 & 0,3–1,0 & \uc{UC40}\\ \hline
Trọng số $R$, $E$, $F$ & 0,6; 0,3; 0,1 & Mỗi trọng số 0–1, tổng bằng 1 & \uc{UC40}\\ \hline
Số món gợi ý tối đa & 10 & 3–30 & \uc{UC40}\\ \hline
Khoảng lập kế hoạch tối đa & 28 ngày & 7–60 ngày & \uc{UC37}\\ \hline
Số lô tối đa trong một thông báo tổng hợp & 5 & 1–10 & \uc{UC31}\\ \hline
\end{longtable}
}

\begin{ucspec}{UC53}{Cấu hình tham số hệ thống}
\ucfield{Tác nhân}{Quản trị viên}
\ucfield{Mô tả}{Cho phép quản trị viên xem và điều chỉnh các tham số vận hành của hệ thống trong Bảng~\ref{tab:tham-so-he-thong}.}
\ucfield{Điều kiện trước}{Quản trị viên đã đăng nhập.}
\ucflow{Luồng thực thi chính}
\ucstep{1}{Quản trị viên}{Mở khu vực Quản trị, chọn “Cấu hình hệ thống”.}
\ucstep{2}{Hệ thống}{Hiển thị các tham số theo nhóm Nhắc nhở, Nhóm gia đình, Bảo mật, Gợi ý món; mỗi tham số kèm giá trị hiện tại, giá trị mặc định và miền hợp lệ.}
\ucstep{3}{Quản trị viên}{Thay đổi một hoặc nhiều tham số và nhấn “Lưu”.}
\ucstep{4}{Hệ thống}{Kiểm tra từng giá trị nằm trong miền hợp lệ và các ràng buộc liên quan, chẳng hạn tổng ba trọng số bằng 1.}
\ucstep{5}{Hệ thống}{Hiển thị bảng tóm tắt các thay đổi (giá trị cũ, giá trị mới) để xác nhận.}
\ucstep{6}{Quản trị viên}{Nhấn “Xác nhận”.}
\ucstep{7}{Hệ thống}{Lưu cấu hình, ghi nhật ký quản trị gồm người thay đổi, thời điểm, giá trị cũ và mới; áp dụng cho các lần xử lý kế tiếp; hiển thị \msg{MSG10}.}
\ucflow{Luồng thực thi mở rộng}
\ucstep{3a}{Quản trị viên}{Chọn “Khôi phục mặc định” cho một tham số: hệ thống điền giá trị mặc định; tiếp tục từ bước 4.}
\ucstep{4a}{Hệ thống}{Có giá trị ngoài miền hợp lệ: hiển thị \msg{MSG27} tại tham số tương ứng; quay lại bước 3.}
\ucstep{6a}{Quản trị viên}{Nhấn “Hủy”: không lưu thay đổi.}
\ucfield{Điều kiện sau}{Bộ tham số mới có hiệu lực; lịch sử thay đổi được lưu để truy vết.}
\ucfield{Quy tắc nghiệp vụ}{\br{BR05}, \br{BR06}, \br{BR12}, \br{BR13}, \br{BR18}, \br{BR19}}
\end{ucspec}

Các use case \uc{UC48}, \uc{UC50}, \uc{UC51} và \uc{UC52} được đặc tả tại Phụ lục~\ref{pl:dac-ta-bo-sung}.

% ---------------------------------------------------------------------
%  2.3. Các yêu cầu phi chức năng
% ---------------------------------------------------------------------
\section{Các yêu cầu phi chức năng}\label{sec:phi-chuc-nang}

Yêu cầu phi chức năng quy định chất lượng hệ thống cần đáp ứng bên cạnh các chức năng đã đặc tả. Các yêu cầu được nhóm theo những đặc tính phù hợp của mô hình chất lượng ISO/IEC 25010~\cite{iso25010}: chức năng, tính dễ dùng, hiệu năng, độ tin cậy, an toàn bảo mật, khả năng bảo trì và khả chuyển; mục này cũng nêu các ràng buộc kỹ thuật. Mỗi yêu cầu có mã \nfr{NFR} và tiêu chí đo lường để kiểm chứng khi kiểm thử. Bảng tổng hợp được trình bày ở cuối mục.

\subsection{Chức năng (Functionality)}\label{ssec:nfr-chuc-nang}

Các yêu cầu trong mục này áp dụng chung cho nhiều use case. Bảng~\ref{tab:dinh-dang-hien-thi} quy định cách hiển thị dữ liệu: số căn phải, chữ căn trái; định dạng số và ngày tháng tuân theo quy ước Việt Nam (\nfr{NFR01}).

{\small
\begin{longtable}{|L{3.0cm}|L{5.0cm}|C{1.8cm}|L{\dimexpr\textwidth-8\tabcolsep-5\arrayrulewidth-3.0cm-5.0cm-1.8cm\relax}|}
\caption{Quy ước định dạng hiển thị chung}\label{tab:dinh-dang-hien-thi}\\
\hline
\hd{Kiểu dữ liệu} & \hd{Định dạng} & \hd{Căn lề} & \hd{Ví dụ}\\ \hline
\endfirsthead
\multicolumn{4}{l}{\footnotesize\itshape Bảng \thetable\ (tiếp theo)}\\ \hline
\hd{Kiểu dữ liệu} & \hd{Định dạng} & \hd{Căn lề} & \hd{Ví dụ}\\ \hline
\endhead
Văn bản & Viết hoa chữ cái đầu câu, cắt bằng dấu ba chấm khi vượt độ rộng & Trái & Thịt ba chỉ\\ \hline
Số lượng & Dấu phẩy thập phân, tối đa hai chữ số thập phân, bỏ số 0 thừa & Phải & 1,5 kg\\ \hline
Tiền tệ & Dấu chấm phân tách hàng nghìn, ký hiệu đồng ở cuối & Phải & 125.000 ₫\\ \hline
Ngày & dd/MM/yyyy; với ngày gần dùng cách nói tương đối & Phải & 14/12/2025; Hôm nay; Ngày mai\\ \hline
Giờ & 24 giờ, HH:mm & Phải & 08:00\\ \hline
Số ngày còn lại & “Còn n ngày”, “Hết hạn hôm nay”, “Quá hạn n ngày” & Phải & Còn 2 ngày\\ \hline
Tuần & Bắt đầu từ thứ Hai & — & T2 08/12 – CN 14/12\\ \hline
\end{longtable}
}

Mẫu báo cáo gợi ý dùng Arial cỡ 14, chữ đen trên nền trắng. Trên thiết bị di động, ứng dụng dùng phông mặc định của hệ điều hành (Roboto trên Android, San Francisco trên iOS), cỡ chữ nội dung tối thiểu 14 sp và tiêu đề từ 18 đến 22 sp. Giao diện chế độ sáng dùng chữ đen hoặc xám đậm trên nền trắng; chế độ tối và thiết lập cỡ chữ của hệ điều hành cũng được hỗ trợ (\nfr{NFR02}). Lỗi nhập liệu hiển thị ngay dưới trường tương ứng, nêu nguyên nhân và cách khắc phục theo danh mục ở Phụ lục~\ref{pl:thong-diep}; thao tác xóa cần xác nhận hoặc có thể hoàn tác trong vài giây (\nfr{NFR03}). Trạng thái hạn dùng được phân biệt bằng màu xanh lá, cam và đỏ, đồng thời luôn có nhãn chữ để người dùng khó phân biệt màu vẫn nhận biết được (\nfr{NFR04}).

\subsection{Tính dễ dùng (Usability)}\label{ssec:nfr-de-dung}

Ứng dụng có thể được sử dụng khi người dùng đang mua sắm hoặc nấu ăn, lúc thao tác không thuận tiện. Vì vậy, giao diện cần dễ thao tác và chỉ rõ vị trí, nguyên nhân, cách khắc phục lỗi nhập liệu. Ba thao tác cốt lõi — thêm mặt hàng, đánh dấu đã mua và thêm thực phẩm vào tủ — cần hoàn thành trong tối đa ba lần chạm từ màn hình chính. Người dùng mới cần tạo được danh sách đầu tiên trong vòng hai phút mà không cần đọc hướng dẫn (\nfr{NFR05}). Để hỗ trợ khả năng tiếp cận, phần tử tương tác có vùng chạm tối thiểu 48 × 48 dp, độ tương phản chữ-nền tối thiểu 4,5 : 1 theo mức AA của WCAG 2.1~\cite{wcag21}, và nhãn cho trình đọc màn hình TalkBack, VoiceOver (\nfr{NFR06}). Giao diện mặc định bằng tiếng Việt, có thể chuyển sang tiếng Anh; chuỗi hiển thị được tách khỏi mã nguồn để thuận tiện bổ sung ngôn ngữ (\nfr{NFR07}).

\subsection{Hiệu năng (Performance)}\label{ssec:nfr-hieu-nang}

Các thao tác chính cần phản hồi nhanh khi người dùng mua sắm. Trên thiết bị tầm trung, thời gian khởi động nguội không quá 3 giây, chuyển màn hình không quá 300 ms và cuộn danh sách 500 phần tử ở mức 60 khung hình mỗi giây (\nfr{NFR08}). Với mạng 4G, ít nhất 95\% yêu cầu API phải phản hồi trong 1 giây; gợi ý món trả kết quả trong 2 giây với 500 công thức và 200 lô thực phẩm; gợi ý tên khi nhập xuất hiện trong 500 ms (\nfr{NFR09}). Trong nhóm gia đình, thay đổi danh sách mua sắm phải đồng bộ tới thiết bị đang trực tuyến trong tối đa 3 giây (\nfr{NFR10}).

\subsection{Độ tin cậy (Reliability)}\label{ssec:nfr-tin-cay}

Do kết nối tại chợ truyền thống và tầng hầm siêu thị có thể không ổn định, ứng dụng cần hỗ trợ ngoại tuyến các chức năng xem, cập nhật danh sách mua sắm, đánh dấu đã mua, xem tủ lạnh và thêm thực phẩm. Thay đổi được lưu vào hàng đợi trên thiết bị và đồng bộ khi có mạng. Khi xảy ra xung đột, thao tác đánh dấu đã mua giữ bản ghi sớm hơn, còn thao tác sửa thuộc tính giữ bản ghi sau cùng; người dùng được thông báo khi thay đổi bị ghi đè (\nfr{NFR11}).

Tác vụ nhắc hạn dùng phải có tính lũy đẳng để chạy lại không tạo thông báo trùng. Mỗi lần gửi thông báo đẩy thất bại được thử lại tối đa ba lần với khoảng chờ tăng dần; ít nhất 99\% thông báo phải được tạo trong vòng 15 phút kể từ khung giờ đã cấu hình (\nfr{NFR12}). Máy chủ cần sẵn sàng tối thiểu 99\% thời gian trong khung 06:00–23:00. Cơ sở dữ liệu được sao lưu hằng ngày, giữ bảy bản gần nhất; các thao tác nhiều bản ghi như \uc{UC21} và \uc{UC41} chạy trong giao dịch nguyên tử (\nfr{NFR13}).

\subsection{An toàn bảo mật (Security)}\label{ssec:nfr-bao-mat}

Ứng dụng lưu thông tin cá nhân và thói quen sinh hoạt của gia đình. Các yêu cầu bảo mật được xây dựng dựa trên OWASP MASVS~\cite{masvs}. Kết nối giữa ứng dụng và máy chủ sử dụng HTTPS với TLS 1.2 trở lên. Mật khẩu được băm bằng thuật toán có muối như bcrypt hoặc Argon2, không lưu trữ hay ghi nhật ký ở dạng rõ (\nfr{NFR14}). Mã truy cập và mã làm mới theo \br{BR03} được lưu trong kho khóa an toàn của hệ điều hành (Keystore trên Android, Keychain trên iOS) và bị thu hồi khi đăng xuất hoặc đổi mật khẩu (\nfr{NFR15}).

Máy chủ kiểm tra quyền của mọi API theo không gian dữ liệu và vai trò người gọi; ẩn nút trên giao diện không được xem là biện pháp phân quyền (\nfr{NFR16}). API xác thực được giới hạn tần suất theo địa chỉ và tài khoản, kết hợp cơ chế khóa tạm tại \br{BR02} và \br{BR03} để hạn chế dò mật khẩu, mã OTP (\nfr{NFR17}). Ứng dụng chỉ xin quyền camera khi người dùng quét mã và chỉ xin quyền thông báo sau khi giải thích mục đích. Người dùng có thể xóa tài khoản cùng dữ liệu cá nhân; hệ thống tuân thủ các quy định hiện hành của Việt Nam về bảo vệ dữ liệu cá nhân (\nfr{NFR18}).

\subsection{Tính dễ bảo trì và khả chuyển (Maintainability, Portability)}\label{ssec:nfr-bao-tri}

Ứng dụng được phát triển bằng một mã nguồn Flutter~\cite{flutter}, hỗ trợ Android 8.0 (API 26) trở lên và iOS 13 trở lên. Giao diện thích ứng với màn hình từ 4,7 đến 6,9 inch và ưu tiên hướng dọc (\nfr{NFR19}). Mã nguồn được chia thành tầng trình bày, nghiệp vụ và dữ liệu. Các quy tắc có tham số được đọc từ cấu hình tại \uc{UC53}; độ bao phủ kiểm thử đơn vị của tầng nghiệp vụ tối thiểu 70\%. API máy chủ có đặc tả OpenAPI và nhật ký có cấu trúc để hỗ trợ truy vết lỗi (\nfr{NFR20}).

\subsection{Các ràng buộc kỹ thuật}\label{ssec:nfr-rang-buoc}

Các ràng buộc kỹ thuật xuất phát từ yêu cầu học phần và các quyết định kiến trúc ban đầu sẽ được cụ thể hóa trong giai đoạn thiết kế. Ứng dụng giao tiếp với máy chủ qua API REST, trao đổi dữ liệu ở định dạng JSON và dùng đường dẫn cơ sở theo quy ước của học phần: \texttt{https://<tên miền>/it4788/}. Các phản hồi dùng chung một bộ mã thống nhất (\nfr{NFR21}). Máy chủ sử dụng cơ sở dữ liệu quan hệ, dịch vụ lưu trữ tệp riêng cho ảnh, Firebase Cloud Messaging~\cite{fcm} để gửi thông báo đẩy trên Android và iOS, cùng bộ lập lịch cho tác vụ định kỳ (\nfr{NFR22}). Thiết bị dùng cơ sở dữ liệu nhúng để lưu đệm khi ngoại tuyến và lưu bản nháp biểu mẫu (\nfr{NFR23}).

Bảng~\ref{tab:tong-hop-nfr} tổng hợp các yêu cầu phi chức năng và tiêu chí đo, làm cơ sở xây dựng ca kiểm thử.

{\small
\begin{longtable}{|C{1.35cm}|L{2.5cm}|L{\dimexpr\textwidth-8\tabcolsep-5\arrayrulewidth-1.35cm-2.5cm-5.0cm\relax}|L{5.0cm}|}
\caption{Tổng hợp các yêu cầu phi chức năng}\label{tab:tong-hop-nfr}\\
\hline
\hd{Mã} & \hd{Đặc tính} & \hd{Yêu cầu} & \hd{Tiêu chí đo}\\ \hline
\endfirsthead
\multicolumn{4}{l}{\footnotesize\itshape Bảng \thetable\ (tiếp theo)}\\ \hline
\hd{Mã} & \hd{Đặc tính} & \hd{Yêu cầu} & \hd{Tiêu chí đo}\\ \hline
\endhead
NFR01 & Chức năng & Định dạng hiển thị thống nhất & Đúng 100\% quy ước của Bảng~\ref{tab:dinh-dang-hien-thi} trên các màn hình\\ \hline
NFR02 & Chức năng & Phông chữ, cỡ chữ, màu nền & Cỡ chữ nội dung $\ge$ 14 sp; hỗ trợ chế độ tối và cỡ chữ hệ thống\\ \hline
NFR03 & Chức năng & Thông báo lỗi tại chỗ, xác nhận khi xóa & Mọi lỗi nhập hiển thị dưới trường; mọi thao tác xóa có xác nhận hoặc hoàn tác\\ \hline
NFR04 & Chức năng & Mã màu trạng thái hạn dùng & Ba màu thống nhất, luôn kèm nhãn chữ\\ \hline
NFR05 & Dễ dùng & Thao tác cốt lõi nhanh & $\le$ 3 lần chạm; người mới tạo danh sách đầu tiên $\le$ 2 phút\\ \hline
NFR06 & Dễ dùng & Khả năng tiếp cận & Vùng chạm $\ge$ 48 dp; tương phản $\ge$ 4,5 : 1; có nhãn cho trình đọc màn hình\\ \hline
NFR07 & Dễ dùng & Đa ngôn ngữ & Tiếng Việt, tiếng Anh; không có chuỗi viết cứng trong mã\\ \hline
NFR08 & Hiệu năng & Hiệu năng giao diện & Khởi động $\le$ 3 s; chuyển màn hình $\le$ 300 ms; cuộn 60 fps với 500 phần tử\\ \hline
NFR09 & Hiệu năng & Thời gian phản hồi máy chủ & 95\% API $\le$ 1 s; gợi ý món $\le$ 2 s; gợi ý tên $\le$ 500 ms\\ \hline
NFR10 & Hiệu năng & Đồng bộ thời gian thực & Thay đổi hiển thị trên thiết bị khác $\le$ 3 s\\ \hline
NFR11 & Tin cậy & Hoạt động ngoại tuyến & Bốn nhóm thao tác chính dùng được khi mất mạng; đồng bộ tự động, không mất dữ liệu\\ \hline
NFR12 & Tin cậy & Độ tin cậy của nhắc nhở & Không gửi trùng; thử lại 3 lần; $\ge$ 99\% thông báo tạo trong 15 phút\\ \hline
NFR13 & Tin cậy & Sẵn sàng, sao lưu, giao dịch & Sẵn sàng $\ge$ 99\% (06:00–23:00); sao lưu hằng ngày, giữ 7 bản\\ \hline
NFR14 & Bảo mật & Mã hóa đường truyền, băm mật khẩu & TLS $\ge$ 1.2; bcrypt hoặc Argon2; không có mật khẩu rõ trong CSDL, nhật ký\\ \hline
NFR15 & Bảo mật & Quản lý phiên & Mã lưu trong Keystore, Keychain; bị thu hồi khi đăng xuất, đổi mật khẩu\\ \hline
NFR16 & Bảo mật & Phân quyền phía máy chủ & 100\% API kiểm tra không gian và vai trò\\ \hline
NFR17 & Bảo mật & Chống dò mật khẩu, OTP & Giới hạn tần suất; khóa tạm theo \br{BR02}, \br{BR03}\\ \hline
NFR18 & Bảo mật & Quyền riêng tư & Xin quyền đúng lúc; hỗ trợ xóa tài khoản và dữ liệu\\ \hline
NFR19 & Khả chuyển & Đa nền tảng & Một mã nguồn Flutter; Android $\ge$ 8.0, iOS $\ge$ 13\\ \hline
NFR20 & Bảo trì & Kiến trúc và chất lượng mã & Phân tầng; tham số hóa quy tắc; bao phủ kiểm thử tầng nghiệp vụ $\ge$ 70\%; có tài liệu OpenAPI\\ \hline
NFR21 & Ràng buộc & Giao tiếp API & REST, JSON, đường dẫn \texttt{/it4788/}, bộ mã phản hồi thống nhất\\ \hline
NFR22 & Ràng buộc & Hạ tầng máy chủ & CSDL quan hệ; lưu trữ tệp riêng; FCM; bộ lập lịch phía máy chủ\\ \hline
NFR23 & Ràng buộc & Lưu trữ cục bộ & CSDL nhúng trên thiết bị cho dữ liệu ngoại tuyến và bản nháp\\ \hline
\end{longtable}
}

Chương~\ref{chap:dac-ta} xác định chín tác nhân và 53 use case trong chín gói chức năng. Chương trình bày biểu đồ use case tổng quan, chín biểu đồ phân rã, đặc tả 24 use case chính cùng các biểu đồ hoạt động và tuần tự cho những luồng phức tạp, và 23 yêu cầu phi chức năng có tiêu chí đo. Phụ lục~\ref{pl:dac-ta-bo-sung} bổ sung đặc tả 29 use case; Phụ lục~\ref{pl:truy-vet} cung cấp ma trận truy vết. Bộ đặc tả bao phủ bảy nhóm chức năng của đề bài và mười chức năng mở rộng, làm cơ sở cho thiết kế giao diện, backend và kiểm thử ở các chương tiếp theo.

